\documentclass[11pt,letterpaper]{article}
\usepackage[T1]{fontenc}
\usepackage[utf8]{inputenc}
\usepackage{amsmath,amsthm,mathtools}
\usepackage{newtxtext,newtxmath}
\usepackage[margin=1in,headheight=15pt]{geometry}
\usepackage{microtype,booktabs,array,longtable,enumitem,xurl}
\usepackage{xcolor,fancyhdr,needspace,tcolorbox}
\definecolor{ink}{HTML}{354E42}
\definecolor{pale}{HTML}{F1F5F2}
\usepackage[colorlinks=true,linkcolor=ink,citecolor=ink,urlcolor=ink]{hyperref}
\hypersetup{pdftitle={Critical-Point Defects in Surreal Arithmetic},
 pdfauthor={Research draft prepared with ChatGPT for Vladimir Reshetnikov},
 pdfsubject={Large-cardinal embeddings, strong Hahn summation, normal-form support, surcomplex numbers, and omnific integers}}
\pagestyle{fancy}\fancyhf{}
\fancyhead[L]{\small\scshape Critical-point defects}
\fancyhead[R]{\small Surreals and large cardinals}
\fancyfoot[C]{\thepage}
\renewcommand{\headrulewidth}{0.3pt}
\setlength{\parskip}{3.5pt}
\setlength{\parindent}{1.15em}
\setlength{\emergencystretch}{3em}
\setlist{leftmargin=1.7em,itemsep=3pt,topsep=5pt}
\clubpenalty=10000\widowpenalty=10000
\allowdisplaybreaks[2]
\numberwithin{equation}{section}
\newtheorem{theorem}{Theorem}[section]
\newtheorem{lemma}[theorem]{Lemma}
\newtheorem{proposition}[theorem]{Proposition}
\newtheorem{corollary}[theorem]{Corollary}
\theoremstyle{definition}
\newtheorem{definition}[theorem]{Definition}
\newtheorem{example}[theorem]{Example}
\newtheorem{question}[theorem]{Question}
\theoremstyle{remark}
\newtheorem{remark}[theorem]{Remark}
\newcommand{\No}{\mathbf{No}}
\newcommand{\On}{\mathbf{On}}
\newcommand{\Oz}{\mathbf{Oz}}
\newcommand{\SC}{\mathbf{SC}}
\newcommand{\R}{\mathbb R}
\newcommand{\C}{\mathbb C}
\newcommand{\Q}{\mathbb Q}
\newcommand{\Z}{\mathbb Z}
\newcommand{\N}{\mathbb N}
\newcommand{\supp}{\operatorname{supp}}
\newcommand{\dom}{\operatorname{dom}}
\newcommand{\ran}{\operatorname{ran}}
\newcommand{\crit}{\operatorname{crit}}
\newcommand{\LT}{\operatorname{LT}}
\newcommand{\ct}{\operatorname{ct}}
\newcommand{\Frac}{\operatorname{Frac}}
\newcommand{\bd}{\operatorname{b}}
\newcommand{\len}{\operatorname{len}}
\newcommand{\Pow}{\mathcal P}
\newcommand{\Sk}{\mathscr S_{<\kappa}}
\newcommand{\Ik}{\mathscr I_{<\kappa}}
\newcommand{\Gam}{\Gamma_j}
\newcommand{\Common}{C_j}
\newcommand{\Hj}{H_j}
\newcommand{\Dj}{D_j}
\newcommand{\Jc}{J_{\C}}
\newcommand{\Hc}{H_{j,\C}}
\newcommand{\Om}{\Omega}
\newcommand{\eps}{\varepsilon}
\newcommand{\imset}[2]{#1\mathbin{``}#2}
\newcommand{\coef}[2]{[\omega^{#1}]\,#2}
\newcommand{\abs}[1]{\lvert#1\rvert}
\newcommand{\set}[1]{\{#1\}}
\newcommand{\res}{\mathbin{\upharpoonright}}
\newcommand{\infideal}{\mathfrak m}
\newcommand{\ZFC}{\mathsf{ZFC}}
\newcommand{\GBC}{\mathsf{GBC}}
\newcommand{\RCF}{\mathsf{RCF}}
\newcommand{\ACF}{\mathsf{ACF}_0}
\newcommand{\repo}{\texttt{0865f043aec1}}
\newcommand{\status}[1]{\begin{tcolorbox}[colback=pale,colframe=ink,boxrule=.45pt,arc=0pt,left=9pt,right=9pt,top=7pt,bottom=7pt]#1\end{tcolorbox}}
\newcommand{\thref}[1]{Theorem~\ref{#1}}
\newcommand{\prref}[1]{Proposition~\ref{#1}}
\newcommand{\leref}[1]{Lemma~\ref{#1}}

\begin{document}
\begin{titlepage}
\thispagestyle{empty}
\vspace*{0.65in}
{\color{ink}\large\scshape Surreal arithmetic / set theory}\par
\vspace{0.35in}
{\huge\bfseries Critical-Point Defects\\[7pt] in Surreal Arithmetic\par}
\vspace{0.27in}
{\Large Large-cardinal embeddings, exact support thresholds,\\[5pt]
 surcomplex fields, and omnific integers\par}
\vspace{0.45in}
{\large A proof-focused research article}\par
\vspace{0.13in}
Prepared for Vladimir Reshetnikov with ChatGPT\par
23 September 2026\par
\vspace{0.4in}
\status{\textbf{Main result.} An elementary embedding $j:V\to M$ with critical
point $\kappa$ induces two ordered-field embeddings $J,H_j:\No\to\No$.
They agree on every real number, ordinal, and Conway monomial, but
\[
 J(x)=H_j(x)\quad\Longleftrightarrow\quad
 \abs{\supp(x)}<\kappa.
\]
Their images have an exactly determined intersection. A single coefficient
of their discrepancy recovers the normal measure derived from $j$.}
\vfill
\textbf{Scope and status.} All central assertions below have written proofs
under explicit hypotheses. The general Hahn and large-cardinal ingredients
are classical; the exact defect and intersection package is offered as a
candidate original contribution. Publication priority has not been certified,
and the article is neither refereed nor Lean-verified. It does not claim a
solution to a named published conjecture or a new consistency proof for a
large-cardinal axiom.
\vspace{0.24in}

\small Repository context: \texttt{VladimirReshetnikov/Surreal}, snapshot
\repo.\par
Standalone source, compiled PDF, proof-status notes, and finite algebraic
regression checks accompany this article.
\end{titlepage}

\pagenumbering{roman}
\begin{abstract}
Large-cardinal embeddings preserve every finitary arithmetic identity, but
need not preserve the external strong sums that define surreal normal forms.
We make this distinction exact. Let $j:V\to M$ be an amenable elementary
embedding into a transitive inner class, with critical point $\kappa$. Its
restriction to surreal sign sequences gives an ordered-field embedding $J$.
The same action on Conway monomials has a unique strongly $\R$-linear
extension $H_j$. If $x$ has normal form indexed by $\theta$, then $J(x)-H_j(x)$
is precisely the subseries of the transformed normal form indexed by
$j(\theta)\setminus\imset j\theta$. Consequently the equalizer of $J$ and
$H_j$ is exactly the real-closed field of surreal numbers with fewer than
$\kappa$ outer normal-form terms. Moreover, their two image fields intersect
exactly in the common image of that equalizer, and are its incomparable
logarithmic and strong-sum closures. For a positive strongly
summable family, preservation by $J$ holds exactly when the index set has
cardinality less than $\kappa$.

The defect has a first term at the first missing index $\kappa$. For
$h_A=\sum_{\alpha\in A}\omega^{-\alpha}$, $A\subseteq\kappa$, its coefficient
at $\omega^{-\kappa}$ is the indicator of membership in the normal measure
derived from $j$. This gives a coefficient realization of classical
ultrafilter evaluation and a precise measurable-cardinal criterion for the
associated Hadamard mask algebra. The field results transfer to surcomplex
numbers and to omnific integers; a purely infinite omnific defect encodes the
same measure while its constant term vanishes. We also obtain a logarithmic
defect homomorphism, an exact exponential compatibility locus, and a
supercompact seed-coefficient construction. The
article distinguishes these contributions from standard Hahn reindexing,
classical derived measures, and the repository's earlier work on birthday
expansions and forcing. Twelve further research questions and a formalization
plan conclude the account.
\end{abstract}
\vspace{4pt}
\noindent\textbf{Keywords.} Surreal number; Conway normal form; strong summation;
measurable cardinal; elementary embedding; omnific integer; surcomplex number;
Hahn field; ultrafilter; supercompactness.
\begingroup
\small
\setlength{\parskip}{0pt}
\tableofcontents
\endgroup
\clearpage
\pagenumbering{arabic}

\section{The questions answered and the novelty boundary}\label{sec:intro}

\subsection{The distinction that drives the construction}
For a set-theoretic elementary embedding, the image $j(I)$ of an infinite set
is generally larger than its pointwise image $\imset jI$. For a surreal
normal form, this is not merely a difference between two encodings of the
same sum. The additional indices can carry additional nonzero monomials.
The resulting discrepancy is visible in arithmetic, even when every original
summand is individually fixed.

This article answers the following concrete question under a measurable-cardinal
hypothesis: \emph{how much of a surreal number is determined by the action of
an embedding on all real coefficients and all monomials?} The answer is
exactly its short-support part. There are two explicitly associated embeddings
that agree on all those generators, and even on every ordinal, but disagree
on every surreal whose outer normal form has at least $\kappa$ terms.
The distinction is a support-cardinality threshold, not a birthday cutoff.

Conway's normal-form theorem and the generalized-series presentation of
$\No$ are the classical background \cite{Conway,Gonshor,vdDE,EhrlichNames}.
Strongly additive maps and monomial reindexing also have a substantial
literature. In particular, Kaplan--Krapp--Serra distinguish strong from
non-strong automorphisms and study their relation to extra surreal structure
\cite{KKS}. The contribution proposed here is not the bare possibility of
non-strong behavior. It is the exact critical-point formula, equalizer,
intersection, and coefficient reconstruction for the specific large-cardinal
construction.

\subsection{Assumption and headline statements}
Throughout the main construction, assume that
\begin{equation}\label{eq:setting}
 j:V\longrightarrow M\subseteq V
\end{equation}
is an elementary class embedding, $M$ is transitive and contains every
ordinal, and $\kappa=\crit(j)$ exists. Class parameters are allowed in the
Separation and Replacement instances used below. A sufficient, familiar
instance is a normal ultrapower embedding from a measurable cardinal.
One can carry out that instance in $\ZFC$ using the measure as a set
parameter; the use of classes is notation for definable constructions.
The general class-embedding version can be read in $\GBC$ with the indicated
embedding as an amenable class. These hypotheses are discussed in
Section~\ref{sec:foundations}.

Let
\begin{equation}\label{eq:shortfield}
 \Sk=\set{x\in\No:\abs{\supp(x)}<\kappa},
\end{equation}
where $\supp(x)$ means the set of \emph{outer exponents} in its Conway normal
form. Zero has empty support. Define $J(x)=j(x)$ on surreal sign sequences,
and define
\begin{equation}\label{eq:Hintro}
 H_j\left(\sum_{\alpha<\theta}r_\alpha\omega^{a_\alpha}\right)
   =\sum_{\alpha<\theta}r_\alpha\omega^{J(a_\alpha)}.
\end{equation}
The central conclusions, proved in Sections~\ref{sec:companion}--\ref{sec:positive},
are
\begin{align}
 \operatorname{Eq}(J,H_j)&=\Sk,\label{eq:eqintro}\\
 J(\No)\cap H_j(\No)&=J(\Sk)=H_j(\Sk),\label{eq:interintro}\\
 J\left(\sum_{i\in I}x_i\right)=\sum_{i\in I}J(x_i)
  &\quad\Longleftrightarrow\quad \abs I<\kappa
   \quad(x_i>0\text{ and strongly summable}).\label{eq:posintro}
\end{align}
The two image fields in \eqref{eq:interintro} are incomparable, not merely
different. Their common image is real closed. Coefficient projection
canonically identifies the two images as immediate valued extensions of
that common field (Theorem~\ref{thm:projection}). An intrinsic version
shows that every positive strong sum of at least $\kappa$ image elements
escapes the elementary image (Theorem~\ref{thm:escape}).

For $A\subseteq\kappa$, put $h_A=\sum_{\alpha\in A}\omega^{-\alpha}$.
One then obtains
\begin{equation}\label{eq:Uintro}
 \coef{-\kappa}{\bigl(J(h_A)-H_j(h_A)\bigr)}
    =\begin{cases}1,&\kappa\in j(A),\\0,&\kappa\notin j(A).
      \end{cases}
\end{equation}
Thus the usual derived measure is recovered from a distinguished coefficient
of the arithmetic defect. The derived measure itself is classical
\cite{Neeman}; the coefficient realization is the new construction being
proposed.

\subsection{What is not being asserted}
We do not establish that arbitrary non-strong surreal embeddings require a
measurable cardinal. We do not prove a new consistency result about
measurability, supercompactness, or elementary self-embeddings of the set
universe. We do not identify the two embeddings as automorphisms. Nor do we
claim that their equality on monomials implies equality on a completed
field: the whole point is that the relevant external sums need not be
preserved by $J$.

The proofs are mathematical arguments, not output from a theorem prover.
The accompanying finite checks test finite identities used in the arguments;
they cannot test an elementary embedding of the universe or a nonprincipal
complete ultrafilter. A targeted literature and repository comparison appears
in Section~\ref{sec:position}; it is not an exhaustive priority search.

\section{Foundations, normal forms, and the meaning of a sum}\label{sec:foundations}

\subsection{Surreal conventions}
We use the sign-sequence realization of $\No$: a surreal is a function from
an ordinal to $\{-,+\}$. Its birthday $\bd(x)$ is the domain of that function.
The numerical order is the first-disagreement order, with
$-<\text{undefined}<+$. Simplicity is proper prefix inclusion.
Ordinals are identified with all-plus sequences.
The distinction between an ordinal and its all-plus sign code will normally
be suppressed, but all applications of $j$ commute with that identification.

The Conway monomial map is denoted by
\[
 \Om(a)=\omega^a.
\]
Here $\omega^a$ is the Conway omega-map, not an abbreviation for
$\exp(a\log\omega)$ with an arbitrarily chosen exponential structure.
It satisfies $\omega^{a+b}=\omega^a\omega^b$, and is strictly increasing.
Every nonzero surreal has a unique normal form
\begin{equation}\label{eq:NF}
 x=\sum_{\alpha<\theta} r_\alpha\omega^{a_\alpha},
 \qquad r_\alpha\in\R\setminus\{0\},\quad
 a_\alpha>a_\beta\quad(\alpha<\beta<\theta),
\end{equation}
where $\theta$ is an ordinal. All supports in this article are sets, even
though the full coefficient and exponent universes are classes.
Write $\len(x)=\theta$ for the length of \eqref{eq:NF} and
$\supp(x)=\{a_\alpha:\alpha<\theta\}$.
For $x\ne0$, define
\[
 \LT(x)=r_0\omega^{a_0},\qquad v(x)=-a_0.
\]
This is the additive natural valuation with values in $(\No,+)$; $v(0)=\infty$.
The infinitesimal ideal of the finite-surreal valuation ring is
\[
 \infideal=\{0\}\cup\{x\in\No:v(x)>0\}.
\]
These conventions use decreasing omega-exponents; sources that use $t^g$
with increasing support have the opposite sign convention for exponents.

\begin{definition}[Strong summability]\label{def:strong}
A set-indexed family $(x_i)_{i\in I}$ is \emph{strongly summable} if
$\bigcup_{i\in I}\supp(x_i)$ is reverse well ordered and each exponent occurs
in the support of only finitely many members. Its strong sum is the unique
normal form obtained by adding those finitely many coefficients at each
exponent and discarding zero coefficients. A map is \emph{strongly
$\R$-linear} if it is $\R$-linear and preserves all such external set-indexed
sums, including their summability.
\end{definition}

The field laws for generalized series, in particular the finiteness of
contributions in a product, are part of the Hahn--Neumann support theorem
\cite{Gonshor,vdDE,KKS}. We use them for sets of exponents. There is no
summation over a proper-class support.

\begin{remark}[Not order-topological convergence]\label{rem:topology}
A displayed strong sum is not defined as the supremum of its partial sums,
and is not asserted to be an order-topological limit in all of $\No$.
For example, the sum $h_\kappa$ constructed below exists by
Definition~\ref{def:strong}. This makes no assertion that its partial sums
converge against every positive surreal neighborhood. Confusing these two
notions would invalidate the threshold argument.
\end{remark}

\subsection{Elementary embeddings and external sets}
The standard large-cardinal facts used here are that the critical point of
\eqref{eq:setting} is an uncountable inaccessible cardinal, that $j$ fixes
$V_\kappa$, and that a measurable cardinal supplies such an embedding by a
normal ultrapower \cite{Neeman}. In particular,
\begin{equation}\label{eq:fixedcoeff}
 j(\R)=\R,\qquad j(r)=r\ (r\in\R),\qquad \R^M=\R^V.
\end{equation}
We take these large-cardinal facts as background, not as results of this
article. When a normal measure $U$ is chosen, the ultrapower is first formed
on set functions with domain $\kappa$ and then collapsed to its transitive
class model. Countable completeness gives well-foundedness. All individual
images and all set-indexed images used below are sets.

We distinguish $j(I)$, an element of $M$, from the external pointwise image
$\imset jI=\{j(i):i\in I\}$, formed in $V$. The latter need not belong to
$M$. Every cardinality without a superscript is computed in $V$. In
particular, $\abs{\supp(x)}<\kappa$ always concerns the original external
support. No cardinal arithmetic in $M$ is silently substituted for it.

\begin{lemma}[The first missing index]\label{lem:index}
For any set $I$,
\[
 j(I)=\imset jI\quad\Longleftrightarrow\quad\abs I<\kappa.
\]
For any ordinal $\theta$, define
$P_\theta=j(\theta)\setminus\imset j\theta$. Then
\[
 P_\theta=\varnothing\quad(\theta<\kappa),\qquad
 \min P_\theta=\kappa\quad(\theta\ge\kappa).
\]
\end{lemma}
\begin{proof}
Choose a bijection $e:\mu\to I$, where $\mu$ is the cardinality of $I$.
If $\mu<\kappa$, then $j(\mu)=\mu$, and every member of $j(I)$ has the
form $j(e)(\alpha)=j(e(\alpha))$ for some $\alpha<\mu$.
Conversely, suppose $\mu\ge\kappa$. The ordinal $\kappa$ belongs to
$j(\mu)$, but is not $j(\alpha)$ for any ordinal $\alpha$: below $\kappa$
all ordinals are fixed, while for $\alpha\ge\kappa$ we have
$j(\alpha)\ge j(\kappa)>\kappa$. Thus $j(e)(\kappa)$ is in $j(I)$ but
not in $\imset jI$, by injectivity of $j(e)$.

For ordinal $\theta<\kappa$, every index is fixed. If $\theta\ge\kappa$,
then $j(\theta)>\kappa$, all smaller indices are in $\imset j\theta$, and
$\kappa$ is not. This proves both assertions.
\end{proof}

\subsection{The absoluteness needed to compare the two fields}
There are two universe changes in play. Applying $j$ moves an input into
$M$. Then the identity inclusion regards the resulting old surreal sign
sequence as a surreal of $V$. We need the standard constructions to agree
under this second passage.

\begin{lemma}[Absolute surreal constructions]\label{lem:absolute}
Let $M\subseteq V$ be transitive models or inner classes of $\ZFC$ with the
same ordinals and the same reals. On inputs that belong to $M$, the identity
inclusion $\No^M\subseteq\No^V$ preserves order, the field operations, the
Conway omega-map, and canonical normal forms. An $M$-set strongly summable
family in $M$ is externally strongly summable in $V$, with the same sum.
\end{lemma}
\begin{proof}
Order and prefix comparison are comparisons of ordinal-length sign
functions, so they are absolute. The simplest separator of a given
\emph{set} cut is constructed by following the forced signs until the
first permissible termination; this construction depends only on its
input sign sequences. Thus the cut construction is absolute. Well-founded
recursion on canonical sign options gives the same negation, addition,
multiplication, and omega-map in both universes. Inverses then agree by
uniqueness in the ordered field.

For normal forms we use the constructive sign-sequence version of the
Conway--Gonshor normal-form theorem \cite[Chapter~5]{Gonshor}, not just the
existence of an unspecified isomorphism with a Hahn field. The conversion
from a reverse well-ordered list of real-coefficient monomials to its
surreal sign sequence is given by the normal-form sign expansion: ordinal
sign concatenations with the prescribed deletions determined by the input
exponent signs and preceding terms. The ordinal operations, restrictions,
and sign comparisons in that construction are absolute on a given set of
inputs. Hence an $M$-normal form represents the same sign sequence in $V$.
Uniqueness of the normal form in $V$ shows that it acquires no extra terms
on passing to $V$.

Finally, a reverse well ordering of an $M$-set has an ordinal order-type
witness in $M$, which remains a witness in $V$; finiteness of an old fiber
is absolute. The old family is therefore still strongly summable. Its
coefficientwise finite sums agree, as do the normal forms just considered.
\end{proof}

\begin{remark}[Where a classical theorem is used]\label{rem:classicalinput}
The normal-form sign-expansion theorem is a classical input to
\leref{lem:absolute}; it is not reproved from the axioms of surreal games
here. This is the one place where an arbitrary Hahn-field identification
would be insufficient: the identification must be the canonical Conway
normal form. A formal verification should first prove this absoluteness
lemma from that canonical construction, rather than assuming that every
isomorphism chosen in $M$ is absolute.
\end{remark}

\begin{proposition}[The elementary-induced surreal embedding]\label{prop:J}
The map $J:\No\to\No$, $J(x)=j(x)$ on sign codes, is an ordered-field
embedding that fixes $\R$, and
\[
 J(\omega^a)=\omega^{J(a)},\qquad
 \bd(J(x))=j(\bd(x)).
\]
It preserves and reflects prefix simplicity. It is elementary in the
ordered-field language.
\end{proposition}
\begin{proof}
Elementary transfer into $M$ preserves the defining surreal constructions;
\leref{lem:absolute} then identifies their values with the values in $V$.
The all-plus coding of ordinals is definable, so it commutes with $j$.
The domain of $j(x)$ is $j(\dom x)$, proving the birthday formula.
Prefix simplicity is a definable statement about sign functions and is
absolute after transfer. Fixing the reals follows from \eqref{eq:fixedcoeff}.
An embedding between real-closed ordered fields is elementary by
quantifier elimination for $\RCF$. Applied formula by formula, the same
argument works for these class fields: it makes no demand for a class-sized
satisfaction set.
\end{proof}

\section{The strong companion}\label{sec:companion}

\subsection{General Hahn reindexing}
We record the standard mechanism in the form needed here. It does not use
large cardinals.

\begin{lemma}[Reindexing exponents]\label{lem:reindex}
Let $f:(\No,+,<)\to(\No,+,<)$ be an additive order embedding. Then
\[
 T_f\left(\sum_a r_a\omega^a\right)=\sum_a r_a\omega^{f(a)}
\]
defines a strongly $\R$-linear ordered-field embedding of $\No$ into itself.
Its image consists exactly of the surreal numbers whose exponent supports
are contained in $f(\No)$. It preserves and reflects the condition of being
an omnific integer.
\end{lemma}
\begin{proof}
An order embedding carries each reverse well-ordered set to a reverse
well-ordered set of the same order type. Thus the displayed normal form
exists and is unique. Coefficientwise addition and scalar multiplication
commute with reindexing. For multiplication, $f(a+b)=f(a)+f(b)$, and
injectivity shows that exactly the same finitely many contributing pairs
occur at corresponding exponents. Hence $T_f(xy)=T_f(x)T_f(y)$. Also
$T_f(1)=1$, so inverses are preserved by the field equation
$T_f(x)T_f(x^{-1})=1$ when $x\ne0$. A nonzero leading coefficient is
unchanged; signs, strict order, and injectivity follow.

For a strongly summable family, the union of its supports and every finite
coefficient fiber are simply reindexed by $f$. This proves strong
$\R$-linearity for arbitrary external set-indexed families.

One inclusion in the image description is immediate. For the other, take a
set support $S\subseteq f(\No)$. Injectivity and class-parameter Replacement
form the set of its unique inverse images. That set is reverse well ordered,
and reindexing its series gives the desired preimage. No assertion is made
that the inverse-image class of an arbitrary class is a set.

An omnific integer has only nonnegative exponents and has integer
coefficient at exponent zero. An additive order embedding fixes zero,
preserves and reflects nonnegativity, and does not change coefficients.
This proves the last assertion.
\end{proof}

The preceding proof is the usual Hahn reindexing argument; compare
\cite{vdDE,KKS} and the exponent-embedding results described in the
repository \cite{Repo}. Its role is to produce a canonical comparison map,
not to supply the critical-point obstruction.

\begin{definition}[Strong companion and defect]\label{def:companion}
For the embedding in \eqref{eq:setting}, let $H_j=T_J$ and set
\[
 D_j(x)=J(x)-H_j(x),\qquad \Gam=J(\No).
\]
The subgroup $\Gam$ is being used as an \emph{exponent group} when it
occurs inside a support condition. Write
\[
 \mathcal H(\Gam)=\{y\in\No:\supp(y)\subseteq\Gam\}.
\]
\end{definition}

\begin{theorem}[Existence and uniqueness of the companion]\label{thm:companion}
The map $H_j$ is a strongly $\R$-linear ordered-field embedding and
\begin{equation}\label{eq:Himage}
 H_j(\No)=\mathcal H(\Gam).
\end{equation}
It is the unique strongly $\R$-linear map that agrees with $J$ on every
Conway monomial. The maps $J$ and $H_j$ agree on every real number and on
every ordinal, as well as on every monomial.
\end{theorem}
\begin{proof}
Apply \leref{lem:reindex} to the additive order embedding $J$.
For any strongly $\R$-linear map with the stated monomial values, applying
strong linearity to the normal form \eqref{eq:NF} forces
\eqref{eq:Hintro}, proving uniqueness.

Agreement on monomials follows from \prref{prop:J}. Real coefficients are
fixed. Every ordinal has a finite Cantor normal form
$\omega^{\beta_0}n_0+\cdots+\omega^{\beta_{m-1}}n_{m-1}$, which is also its
Conway normal form. Finite additivity therefore gives agreement on every
ordinal, irrespective of its size.
\end{proof}

\begin{remark}[A crucial noncommutation]\label{rem:omega}
The formula for the companion is
\[
 H_j(\omega^a)=\omega^{J(a)},
\]
not $\omega^{H_j(a)}$ in general. In the former expression, $J$ acts on an
\emph{exponent}; in $H_j(x)$, it reindexes the outer exponents of $x$.
Failing to distinguish these two levels would incorrectly make the two
embeddings identical.
\end{remark}

\section{The exact normal-form defect}\label{sec:defect}

\subsection{The transformed index set}
Fix a nonzero $x$ with normal form \eqref{eq:NF}. Regard the coefficient
list $r$ and the exponent list $a$ as set functions with domain $\theta$.
Put
\[
 r^*=j(r),\qquad a^*=j(a),\qquad
 P_\theta=j(\theta)\setminus\imset j\theta.
\]
The transformed lists have domain $j(\theta)$. By elementary transfer and
\leref{lem:absolute}, the full \emph{external} normal form of $J(x)$ is
\begin{equation}\label{eq:Jnormal}
 J(x)=\sum_{\beta<j(\theta)}r^*_\beta\omega^{a^*_\beta}.
\end{equation}
For every $\alpha<\theta$,
\begin{equation}\label{eq:oldindices}
 r^*_{j(\alpha)}=r_\alpha,\qquad
 a^*_{j(\alpha)}=J(a_\alpha).
\end{equation}
The extra indices in $P_\theta$ will be called \emph{missing-image indices}.
They are ordinary set indices in $V$, not informal nonstandard positions.

\begin{theorem}[Exact defect and first term]\label{thm:defect}
For $x$ as above,
\begin{equation}\label{eq:defectformula}
 D_j(x)=\sum_{\beta\in P_\theta}r^*_\beta\omega^{a^*_\beta}.
\end{equation}
Every exponent occurring on the right lies outside $\Gam=J(\No)$.
Consequently
\begin{equation}\label{eq:exactequalizer}
 D_j(x)=0\quad\Longleftrightarrow\quad\theta<\kappa
 \quad\Longleftrightarrow\quad\abs{\supp(x)}<\kappa.
\end{equation}
If $\theta\ge\kappa$, then
\begin{equation}\label{eq:firstterm}
 \LT(D_j(x))=r^*_{\kappa}\omega^{a^*_{\kappa}},\qquad
 v(D_j(x))=-a^*_{\kappa}.
\end{equation}
The same statements include $x=0$ by giving it empty support.
\end{theorem}
\begin{proof}
The normal form \eqref{eq:Jnormal} is a strongly summable family. Split it
into the subfamilies on $\imset j\theta$ and $P_\theta$. Subfamilies of a
strongly summable family remain strongly summable. By
\eqref{eq:oldindices}, the first subfamily sums to $H_j(x)$.
Subtracting it leaves exactly \eqref{eq:defectformula}.

Suppose, toward a contradiction, that $a^*_\beta=J(b)$ for some
$\beta\in P_\theta$ and $b\in\No$. Then in $M$,
$j(b)\in\ran(j(a))$. Elementarity gives $b\in\ran(a)$ in $V$, so
$b=a_\alpha$ for some $\alpha<\theta$. But then
$a^*_\beta=J(a_\alpha)=a^*_{j(\alpha)}$. Since $a^*$ is strictly decreasing,
it is injective, and $\beta=j(\alpha)$, a contradiction. Thus all the
missing-image exponents lie outside $\Gam$.

All $r^*_\beta$ are nonzero, and all $a^*_\beta$ are distinct. There is
therefore no cancellation between the terms of \eqref{eq:defectformula}.
The defect vanishes precisely when $P_\theta$ is empty, which by
\leref{lem:index} is precisely when $\theta<\kappa$.
Since $\kappa$ is an initial ordinal, the latter condition is equivalent to
$\abs\theta<\kappa$. The bijection $\alpha\mapsto a_\alpha$ converts it
into the last condition of \eqref{eq:exactequalizer}.

For $\theta\ge\kappa$, the least missing-image index is $\kappa$.
The exponent list is strictly decreasing, so its term at that index is
the leading term of the defect, proving \eqref{eq:firstterm}.
\end{proof}

\begin{corollary}[Same leading scale, different tails]\label{cor:leading}
For every $x\ne0$,
\[
 \LT(J(x))=\LT(H_j(x))=r_0\omega^{J(a_0)},\qquad
 v(J(x))=v(H_j(x))=J(v(x)).
\]
If $D_j(x)\ne0$, its valuation is strictly larger than this common value.
\end{corollary}
\begin{proof}
Index zero is fixed and is never a missing-image index. The first missing
index, when present, is later in a strictly decreasing exponent list.
\end{proof}

\begin{corollary}[Exact omega-map compatibility locus]\label{cor:omegalocus}
For every $x\in\No$,
\[
 H_j(\omega^x)=\omega^{H_j(x)}
   \quad\Longleftrightarrow\quad x\in\Sk.
\]
Thus $J$ preserves the omega-map globally but is not strongly additive,
whereas $H_j$ is strongly additive but does not preserve the omega-map
globally.
\end{corollary}
\begin{proof}
By \prref{prop:J} and the companion formula,
$H_j(\omega^x)=\omega^{J(x)}$. The omega-map is injective, so the desired
equality is equivalent to $J(x)=H_j(x)$. Apply \thref{thm:defect}.
The existence of a long-support witness is established explicitly in
Section~\ref{sec:witness}.
\end{proof}

\subsection{Why small support is not small birthday}
Every monomial belongs to $\Sk$, however large its exponent or birthday.
For example, $\omega^b\in\Sk$ for every surreal $b$. This excludes any
interpretation of \eqref{eq:exactequalizer} as agreement only on a
set-sized initial segment of $\No$.
Conversely, the simplest witness below has $\kappa$ outer terms and falls
outside the equalizer even though all its individual monomials are fixed
by $J$. The cardinal bound applies to the number of terms, not to the
size of a single exponent.

The distinction from the classical bounded birthday fields $\No(\lambda)$
is substantial. Van den Dries--Ehrlich analyze those fields and their
restricted analytic and exponential closures \cite[Proposition~4.7 and
Corollary~5.5]{vdDE}. We are not renaming those set-sized fields; $\Sk$ is a
proper class because it contains every monomial.

\section{The equalizer field and the exact image intersection}\label{sec:intersection}

\begin{theorem}[Equalizer and common image]\label{thm:intersection}
Let $\Common=J(\Sk)=H_j(\Sk)$. Then
\begin{align}
 \operatorname{Eq}(J,H_j)&=\Sk,\label{eq:equalizer}\\
 J(\No)\cap H_j(\No)&=\Common,\label{eq:intersection}\\
 \Common&=\{y\in\No:\abs{\supp(y)}<\kappa
                   \text{ and }\supp(y)\subseteq\Gam\}.\label{eq:Cdescription}
\end{align}
Both $\Sk$ and $\Common$ are real-closed subfields of $\No$.
\end{theorem}
\begin{proof}
The equalizer assertion is \thref{thm:defect}. Suppose
$y=J(x)\in H_j(\No)$. The image description \eqref{eq:Himage} says that
all exponents in the normal form of $y$ lie in $\Gam$. If $x$ had at
least $\kappa$ normal-form terms, \thref{thm:defect} would exhibit a
nonzero term of $J(x)$ with exponent outside $\Gam$. Hence $x\in\Sk$,
and $y=J(x)=H_j(x)\in\Common$. The opposite inclusion is immediate.
Reindexing preserves the cardinality of support, and the inverse-support
construction in \leref{lem:reindex} proves \eqref{eq:Cdescription}.

The equalizer of two field embeddings is a subfield, so $\Sk$ is a
subfield. To prove real closedness without relying on a separate
bounded-support theorem, let $0\ne p\in\Sk[X]$. If $r_1<\cdots<r_m$ are all
its distinct real roots in $\No$, then $J(r_1)<\cdots<J(r_m)$ and
$H_j(r_1)<\cdots<H_j(r_m)$ are both the full ordered list of real roots
of the same transformed polynomial. Indeed, the images of $\No$ under
these embeddings are real closed, and an ordered extension of a real-closed
field cannot add a new real algebraic element. Each ordered position
therefore agrees: $J(r_i)=H_j(r_i)$. Every such root lies in $\Sk$.

In particular, each positive member of $\Sk$ has its positive square root
in $\Sk$, and each odd-degree polynomial over $\Sk$ has a root there.
This characterizes real closedness. The isomorphic common image
$\Common$ is real closed as well.
\end{proof}

\begin{remark}[Domain equalizer versus image intersection]\label{rem:twocommon}
The equalizer is $\Sk$, whereas the intersection is its \emph{image}
$\Common$. The two are not being identified. An arbitrary short-support
series can have exponents outside $J(\No)$, and then it is not in
$\Common$. In \eqref{eq:Cdescription}, both the support bound and the
exponent-image condition are essential.
\end{remark}

\begin{corollary}[Transcendental separation]\label{cor:transcendence}
Every element of $J(\No)\setminus\Common$ is transcendental over
$H_j(\No)$, and every element of $H_j(\No)\setminus\Common$ is
transcendental over $J(\No)$. In particular, each is transcendental over
$\Common$.
\end{corollary}
\begin{proof}
Both image fields are real closed. An element of one outside their
intersection is not in the other. A real-closed field has no proper
algebraic ordered extension, so that element cannot be algebraic over
the other field.
\end{proof}

\subsection{A canonical isomorphism by erasing nonimage exponents}
Define a coefficient projection on all surreals by
\begin{equation}\label{eq:projection}
 \Pi_j(y)=\sum_{a\in\supp(y)\cap\Gam}
                  (\coef{a}{y})\,\omega^a.
\end{equation}
This retains a subseries; it is not generally an initial normal-form
truncation. Its support is a set, formed by Separation with the embedding
parameter.

\begin{theorem}[Isomorphic immediate image fields]\label{thm:projection}
The projection $\Pi_j$ is strongly $\R$-linear and idempotent, with image
$H_j(\No)$, and
\begin{equation}\label{eq:projectionJ}
 \Pi_j\circ J=H_j.
\end{equation}
Its restriction is an ordered-field isomorphism
\begin{equation}\label{eq:fieldiso}
 \Pi_j\res J(\No):J(\No)\xrightarrow{\ \cong\ }H_j(\No)
\end{equation}
that fixes $\Common$ pointwise and fixes every Conway monomial in either
image. Both image fields are immediate valued extensions of $\Common$:
they have the same value group $\Gam$ and residue field $\R$.
The projection is not a field homomorphism on all of $\No$.
\end{theorem}
\begin{proof}
Taking a subseries is compatible with coefficientwise finite sums, and
commutes with an arbitrary strong sum: one either retains or deletes all
contributions at each exponent. This proves strong $\R$-linearity.
The image and idempotence follow from \eqref{eq:Himage}.
The defect theorem states that $J(x)$ consists of $H_j(x)$ supported on
$\Gam$, plus a series supported outside $\Gam$. Therefore
$\Pi_j(J(x))=H_j(x)$.

The maps $J$ and $H_j$ are ordered-field isomorphisms onto their respective
images. Hence \eqref{eq:projectionJ} identifies the restriction in
\eqref{eq:fieldiso} with $H_j\circ J^{-1}$, which is an ordered-field
isomorphism. It fixes every element of the intersection, since those
supports are contained in $\Gam$. The monomials in either image are
exactly $\{\omega^{J(a)}:a\in\No\}$, and they already lie in $\Common$.
For the assertion about the monomials of $J(\No)$, a monomial image must
have normal-form length one and leading coefficient one, so its preimage
also does.

Every nonzero member of either image has valuation in $\Gam$, and the
monomials just listed realize every value in $\Gam$. They belong to
$\Common$, so its value group is the same. All three fields contain the
constant copy of $\R$ and have residue field $\R$. Equality of these
value groups and residue fields is precisely immediacy; the value groups
here are allowed to be classes.

Finally choose $a\notin\Gam$, for instance $a=\kappa$. Then also
$-a\notin\Gam$, and
\[
 \Pi_j(\omega^a)=\Pi_j(\omega^{-a})=0,
 \qquad \Pi_j(\omega^a\omega^{-a})=\Pi_j(1)=1.
\]
Thus global multiplicativity fails. The fact that multiplicativity holds
on $J(\No)$ is a genuine restriction on the tails of that image field.
\end{proof}

This gives transcendence, not algebraic independence between the two image
fields. Exact intersection alone does not establish linear disjointness;
that stronger question is left open in Section~\ref{sec:questions}.

\section{Explicit witnesses and incomparability}\label{sec:witness}

For every ordinal $\eta$, define the genuine strong sum
\begin{equation}\label{eq:heta}
 h_\eta=\sum_{\alpha<\eta}\omega^{-\alpha}.
\end{equation}
The set $\{-\alpha:\alpha<\eta\}$ is strictly decreasing in its displayed
order, and each exponent occurs once, so this is a legitimate normal form.
We always use surreal negation and subtraction in exponents; they are not
ordinal subtraction.

\begin{theorem}[The basic witness]\label{thm:witness}
For $h_\kappa$ one has
\begin{align}
 H_j(h_\kappa)&=h_\kappa,\label{eq:Hh}\\
 J(h_\kappa)&=h_{j(\kappa)},\label{eq:Jh}\\
 D_j(h_\kappa)&=
  \sum_{\kappa\le\beta<j(\kappa)}\omega^{-\beta}
  =\omega^{-\kappa}(1+\eps),\qquad 0<\eps\in\infideal.
  \label{eq:Dh}
\end{align}
Moreover,
\[
 h_\kappa\in H_j(\No)\setminus J(\No),\qquad
 h_{j(\kappa)}\in J(\No)\setminus H_j(\No).
\]
Thus $J(\No)$ and $H_j(\No)$ are incomparable real-closed subfields of
$\No$, although the embeddings agree on all reals, ordinals, and monomials.
\end{theorem}
\begin{proof}
Each $\alpha<\kappa$ is fixed, so $J(-\alpha)=-\alpha$, and
\eqref{eq:Hh} follows. The definition of the whole list in
\eqref{eq:heta} transfers to the list of length $j(\kappa)$, giving
\eqref{eq:Jh}. Their difference is \eqref{eq:Dh}. Factoring out its first
monomial leaves 1 plus a nonempty series with strictly negative exponents
and positive coefficients. That remainder is a positive infinitesimal.
It is nonempty because $j(\kappa)$ is a limit cardinal strictly above
$\kappa$.

If $h_\kappa=J(x)$, then the length of its normal form is
$j(\len(x))=\kappa$, contradicting that $\kappa$ is not in the ordinal
range of $j$. Hence $h_\kappa$ is not in $J(\No)$.
On the other hand, the normal form of $h_{j(\kappa)}$ contains the exponent
$-\kappa$. This exponent is not $J(b)$ for any surreal $b$: if it were,
$J(-b)=\kappa$, and elementarity would make $-b$ an ordinal, again
contradicting the missing-index fact. The image description for $H_j$
therefore excludes $h_{j(\kappa)}$ from its image.
\end{proof}

\begin{example}[A short-support element outside the common image]
By \prref{prop:J}, $J(\kappa)=j(\kappa)>\kappa$. Consequently
$\omega^{-\kappa}$ is not in $\Common$: its exponent is outside $\Gam$.
This is a one-term surreal, so it nevertheless belongs to $\Sk$.
It concretely illustrates Remark~\ref{rem:twocommon}. The two assertions
``$x$ has short support'' and ``$x$ lies in the common image'' really are
different assertions.
\end{example}

\section{The exact strong-summation threshold}\label{sec:positive}

The previous results compare the images of a \emph{single normal form}.
There is a separate, equally sharp statement about an arbitrary family
whose individual members may themselves have very long supports.

\begin{theorem}[Positive-family threshold]\label{thm:positive}
Let $F:I\to\No$ be an external set-indexed strongly summable family.
Then $(J(F(i)))_{i\in I}$ is strongly summable. If $\abs I<\kappa$, then
\begin{equation}\label{eq:smallfamily}
 J\left(\sum_{i\in I}F(i)\right)=\sum_{i\in I}J(F(i)).
\end{equation}
If in addition $F(i)>0$ for every $i\in I$, then
\eqref{eq:smallfamily} holds \emph{if and only if} $\abs I<\kappa$.
For $\abs I\ge\kappa$ in this positive case,
\begin{equation}\label{eq:positivegap}
 J\left(\sum_{i\in I}F(i)\right)-\sum_{i\in I}J(F(i))>0.
\end{equation}
\end{theorem}
\begin{proof}
The family $j(F)$ is strongly summable over $j(I)$ in $M$, and hence
externally in $V$ by \leref{lem:absolute}. Its sum is
$J(\sum_{i\in I}F(i))$. Restrict it to the old-index part $\imset jI$:
its values there are $j(F)(j(i))=J(F(i))$. This proves external
summability of the termwise image family. Splitting the full family yields
\begin{equation}\label{eq:familydefect}
 J\left(\sum_{i\in I}F(i)\right)-\sum_{i\in I}J(F(i))
    =\sum_{u\in j(I)\setminus\imset jI}j(F)(u).
\end{equation}
If $\abs I<\kappa$, the right side is empty by \leref{lem:index}.

Suppose now that every $F(i)>0$. Transfer gives $j(F)(u)>0$ for all
$u\in j(I)$. A nonempty strongly summable family of positive surreals has
positive sum: take the largest exponent in the union of its supports.
Every term contributing at that exponent has it as a leading exponent,
so its coefficient is positive; only finitely many contribute, and their
sum is positive. When $\abs I\ge\kappa$, the missing-index part is
nonempty by \leref{lem:index}. Applying this observation to
\eqref{eq:familydefect} proves \eqref{eq:positivegap}.
\end{proof}

\begin{remark}[Why the positivity hypothesis is necessary]
Let $I=\kappa\times\{0,1\}$ and put
$F(\alpha,0)=\omega^{-\alpha}$,
$F(\alpha,1)=-\omega^{-\alpha}$. This is strongly summable and has sum
zero. The termwise image also sums to zero, so preservation holds even
though $\abs I=\kappa$. One cannot replace the positive-family conclusion
by an unconditional cardinality criterion for signed families.
The exact criterion in general is that the missing-index sum in
\eqref{eq:familydefect} vanish.
\end{remark}

\begin{corollary}[First failure cardinal]\label{cor:firstfailure}
The least cardinality of a positive strongly summable family on which $J$
fails to commute with strong summation is exactly $\kappa$.
\end{corollary}
\begin{proof}
All smaller families are preserved, and the monomial family defining
$h_\kappa$ is a positive witness of size $\kappa$.
\end{proof}

There is no conflict between this theorem and the defect of an individual
long-support $x$. A singleton family $\{x\}$ is always preserved by $J$.
Its normal-form decomposition may have $\kappa$ or more terms, and that
is the different family on which $J$ and $H_j$ disagree.

\subsection{An intrinsic closure theorem for the image}
The threshold can be stated without first specifying a family in the
domain of $J$. This strengthens the noncommutation result to an actual
failure of membership in the image.

\begin{theorem}[Strong-sum closure and positive escape]\label{thm:escape}
Let $(y_i)_{i\in I}$ be an external strongly summable family in $J(\No)$,
and let $x_i$ be the unique preimage with $J(x_i)=y_i$. Then $(x_i)$ is
strongly summable and
\begin{equation}\label{eq:membershipsum}
 \sum_{i\in I}y_i\in J(\No)
 \quad\Longleftrightarrow\quad
 \sum_{i\in I}y_i=J\left(\sum_{i\in I}x_i\right).
\end{equation}
In particular, $J(\No)$ is closed under all external strongly summable
families of cardinality less than $\kappa$. If every $y_i>0$, then
\begin{equation}\label{eq:positiveescape}
 \sum_{i\in I}y_i\in J(\No)
   \quad\Longleftrightarrow\quad\abs I<\kappa.
\end{equation}
In contrast, $H_j(\No)$ is closed under every external set-indexed strong
sum of its members.
\end{theorem}
\begin{proof}
Replacement with the embedding parameter gives the set-indexed preimage
family. For every $a$, the normal-form defect theorem implies
\[
 \coef{J(a)}{y_i}=\coef{a}{x_i}.
\]
Thus $\imset J{\bigcup_i\supp(x_i)}$ is a subset of the reverse
well-ordered union $\bigcup_i\supp(y_i)$. Pulling this suborder back by
the order embedding $J$ proves reverse well ordering of
$\bigcup_i\supp(x_i)$. An exponent $a$ occurs in $x_i$ exactly when
$J(a)$ occurs in $y_i$, so the coefficient fibers are finite as well.
Hence the preimage family is strongly summable.

Let $x=\sum_i x_i$ and $y=\sum_i y_i$. At every exponent $J(a)$ the
finite coefficient sums give
\[
 \coef{J(a)}{y}=\sum_i\coef{a}{x_i}
        =\coef{a}{x}=\coef{J(a)}{J(x)}.
\]
Therefore $J(x)-y$ has support disjoint from $\Gam$.
But every nonzero element of $J(\No)$ has its leading exponent in
$\Gam$, by Corollary~\ref{cor:leading}. If $y\in J(\No)$, the
difference $J(x)-y$ is in that field and must consequently be zero.
This proves the forward implication in \eqref{eq:membershipsum}; the
reverse implication is immediate.

For $\abs I<\kappa$, Theorem~\ref{thm:positive} gives equality, without
requiring positive terms. If all $y_i>0$, then all $x_i>0$, and that
theorem makes equality equivalent to $\abs I<\kappa$. Finally, a strong
sum of series supported on $\Gam$ is again supported on $\Gam$;
\eqref{eq:Himage} proves the closure assertion for $H_j(\No)$.
\end{proof}

\section{Recovering a normal measure from one coefficient}\label{sec:measure}

\subsection{Boolean masks at ordinal exponents}
For each $A\subseteq\kappa$, define
\begin{equation}\label{eq:maskA}
 h_A=\sum_{\alpha\in A}\omega^{-\alpha}.
\end{equation}
The notation is unambiguous: the index set is ordered by the ordinal order,
and zeros are omitted from the full $\kappa$-list. In particular,
$h_\varnothing=0$ and $h_\kappa$ agrees with \eqref{eq:heta}.

\begin{theorem}[The measure-coefficient formula]\label{thm:measure}
For every $A\subseteq\kappa$,
\begin{align}
 H_j(h_A)&=h_A,\label{eq:maskfixed}\\
 J(h_A)&=\sum_{\beta\in j(A)}\omega^{-\beta},\label{eq:maskJ}\\
 \coef{-\kappa}{D_j(h_A)}&=\mathbf 1_{\{\kappa\in j(A)\}}.
 \label{eq:measurecoef}
\end{align}
Consequently
\begin{equation}\label{eq:derivedU}
 U_j=\set{A\subseteq\kappa:\coef{-\kappa}{D_j(h_A)}=1}
\end{equation}
is the nonprincipal normal $\kappa$-complete ultrafilter on $\kappa$
derived from $j$.
\end{theorem}
\begin{proof}
The original exponents $-\alpha$, $\alpha<\kappa$, are fixed by $J$, so
\eqref{eq:maskfixed} follows from the companion formula. Transferring the
whole mask construction gives \eqref{eq:maskJ}. Since $\kappa$ is not an
original index, $h_A$ has no coefficient at $\omega^{-\kappa}$.
The coefficient in $J(h_A)$ is 1 precisely when $\kappa\in j(A)$.
This proves \eqref{eq:measurecoef} and identifies \eqref{eq:derivedU}
with the classical derived measure.

For completeness, its properties can be read directly from this last
identification. Since $\kappa\in j(\kappa)$, the whole space is in $U_j$
and the empty set is not. Elementary transfer of complements and finite
intersections makes it an ultrafilter. For $\alpha<\kappa$,
$j(\{\alpha\})=\{\alpha\}$ does not contain $\kappa$, so it is
nonprincipal. If $\mu<\kappa$ and $A_\xi\in U_j$ for $\xi<\mu$, then
$j(\mu)=\mu$ and
\[
 \kappa\in\bigcap_{\xi<\mu}j(A_\xi)
       =j\left(\bigcap_{\xi<\mu}A_\xi\right),
\]
which proves $\kappa$-completeness.

Finally, suppose $A\in U_j$ and $f:A\to\kappa$ is regressive, meaning
$f(\alpha)<\alpha$ for all nonzero $\alpha\in A$. Then
$j(f)(\kappa)<\kappa$, so $j(f)(\kappa)=\eta=j(\eta)$ for some
$\eta<\kappa$. Therefore
\[
 \kappa\in j\bigl(\{\alpha\in A:f(\alpha)=\eta\}\bigr).
\]
The indicated fiber is in $U_j$, proving normality.
\end{proof}

The ultrafilter properties just proved are standard properties of derived
measures \cite{Neeman}. The new information is their simultaneous realization
by \eqref{eq:measurecoef} within one explicitly given surreal coefficient
space. No real-valued probability measure on all surreal sets is being
introduced.

\begin{corollary}[Distinguishing normal measures]\label{cor:distinctU}
Suppose $U$ and $W$ are distinct normal measures on the same $\kappa$, with
normal ultrapower embeddings $j_U$ and $j_W$. Their coefficient functionals
\[
 A\longmapsto\coef{-\kappa}{D_{j_U}(h_A)},\qquad
 A\longmapsto\coef{-\kappa}{D_{j_W}(h_A)}
\]
are distinct. Thus the first failure cardinal does not determine the
coefficient-level defect data.
\end{corollary}
\begin{proof}
Choose $A$ in exactly one of $U,W$. In the normal ultrapower construction,
the measure derived at the seed $\kappa$ is the original measure.
Apply \thref{thm:measure}.
\end{proof}

\begin{remark}[A zero probe does not mean a zero defect]
If $A\subseteq\kappa$ has cardinality $\kappa$ but does not belong to
$U_j$, then $D_j(h_A)\ne0$ by the exact support theorem, although its
coefficient at $\omega^{-\kappa}$ is zero. Such sets exist: partition
$\kappa$ into two sets of size $\kappa$, at most one of which is in the
ultrafilter. The first missing \emph{normal-form index} is always
$\kappa$ for a long support, but its associated \emph{exponent} depends
on the exponent list. The fixed mask probe need not be the leading
nonzero term of every mask defect.
\end{remark}

\subsection{Real-valued masks and the Hadamard product}
For $f:\kappa\to\R$, set
\begin{equation}\label{eq:hf}
 h_f=\sum_{\alpha<\kappa} f(\alpha)\omega^{-\alpha}.
\end{equation}
This includes zero coefficients with the convention that they are suppressed
in the canonical normal form. Let $\mathcal B_\kappa$ denote the set of
all these masks. Give it its coefficientwise, or Hadamard, product
\begin{equation}\label{eq:hadamard}
 h_f\odot h_g=h_{fg},\qquad (fg)(\alpha)=f(\alpha)g(\alpha).
\end{equation}
Together with ordinary addition and real scalar multiplication, this makes
$\mathcal B_\kappa$ an $\R$-algebra isomorphic to $\R^\kappa$.
Its multiplicative unit is $h_\kappa=h_{\mathbf 1}$, \emph{not} the surreal
number 1. The map from scalars into this algebra is
$r\mapsto r h_\kappa$.

\begin{theorem}[Ultrafilter evaluation as a defect coefficient]\label{thm:character}
The functional
\begin{equation}\label{eq:chi}
 \chi_j(h_f)=\coef{-\kappa}{\bigl(J(h_f)-H_j(h_f)\bigr)}
\end{equation}
satisfies
\begin{equation}\label{eq:chievaluation}
 \chi_j(h_f)=j(f)(\kappa).
\end{equation}
It is a unital $\R$-algebra homomorphism
$(\mathcal B_\kappa,+,\odot)\to\R$. Its value is the unique $r\in\R$
for which $f^{-1}(\{r\})\in U_j$.
\end{theorem}
\begin{proof}
All original exponents are fixed, so $H_j(h_f)=h_f$. The whole transformed
mask is
\[
 J(h_f)=\sum_{\beta<j(\kappa)}j(f)(\beta)\omega^{-\beta}.
\]
Reading its coefficient at $-\kappa$ proves
\eqref{eq:chievaluation}. The value lies in $\R$, since $j(f)$ takes
values in $j(\R)=\R$. If $r=j(f)(\kappa)$, then $j(r)=r$ and
\[
 \kappa\in j\bigl(f^{-1}(\{r\})\bigr),
\]
which gives the asserted fiber. Two distinct fibers cannot both belong to
an ultrafilter. Transferring pointwise addition and multiplication gives
\[
 j(f+g)(\kappa)=j(f)(\kappa)+j(g)(\kappa),\qquad
 j(fg)(\kappa)=j(f)(\kappa)j(g)(\kappa).
\]
Constant functions take their constant real value. These are exactly the
claimed algebra-homomorphism laws under \eqref{eq:hadamard}.
\end{proof}

\status{\textbf{Product warning.} The character theorem concerns the Hadamard
product $\odot$. It does not say that coefficient extraction at
$\omega^{-\kappa}$ is multiplicative for ordinary surreal multiplication.
Ordinary multiplication convolves exponents; Hadamard multiplication
multiplies coefficients at matching exponents.}

Infinite direct-product homomorphisms and their ultrafilter factorizations
are a classical subject; see, for example, Bergman \cite{Bergman}.
The next result records an exact large-cardinal criterion in this particular
mask presentation. Its underlying ultrafilter characterization is not claimed
to be new.

\begin{theorem}[A mask criterion for measurability]\label{thm:measurablecriterion}
Let $\kappa$ be an uncountable cardinal, without initially assuming an
embedding $j$. The following are equivalent.
\begin{enumerate}[label=\textup{(\roman*)}]
\item $\kappa$ is measurable.
\item There is a unital $\R$-algebra homomorphism
$\chi:(\mathcal B_\kappa,+,\odot)\to\R$ such that
$\chi(h_{\{\alpha\}})=0$ for every $\alpha<\kappa$, and for every
pairwise disjoint family $(A_\xi)_{\xi<\mu}$ with $\mu<\kappa$,
\begin{equation}\label{eq:smallidempotent}
 \chi\left(h_{\bigcup_{\xi<\mu}A_\xi}\right)
    =\sum_{\xi<\mu}\chi(h_{A_\xi}).
\end{equation}
\end{enumerate}
The scalar sum in \eqref{eq:smallidempotent} has at most one nonzero term,
so no infinite real summation convention is needed. Under measurability,
$\chi$ can be chosen to be the defect functional \eqref{eq:chi} of a
normal ultrapower embedding.
\end{theorem}
\begin{proof}
Assume (i) and choose a normal measure and its ultrapower. The functional
from \thref{thm:character} vanishes on singleton masks. Its value on each
indicator mask is 0 or 1 according to membership in the measure.
For fewer than $\kappa$ disjoint sets, their union belongs to a
$\kappa$-complete ultrafilter exactly when one of them belongs to it:
otherwise all their complements have intersection in the ultrafilter.
This proves \eqref{eq:smallidempotent}.

Conversely, assume (ii). Each $h_A$ is a Hadamard idempotent, so
$\chi(h_A)^2=\chi(h_A)$ and hence $\chi(h_A)\in\{0,1\}$.
Define $U=\{A\subseteq\kappa:\chi(h_A)=1\}$. The identities
\[
 h_A\odot h_B=h_{A\cap B},\qquad
 h_{\kappa\setminus A}=h_\kappa-h_A
\]
show that $U$ is an ultrafilter. For example, upward closure follows from
$h_A\odot h_B=h_A$ when $A\subseteq B$, and exactly one of a set and
its complement has value 1. The singleton hypothesis makes it
nonprincipal.

To verify $\kappa$-completeness, let $A_\xi\in U$ for $\xi<\mu<\kappa$.
Set $B_\xi=\kappa\setminus A_\xi$, and disjointify the family by
\[
 C_\xi=B_\xi\setminus\bigcup_{\eta<\xi}B_\eta.
\]
Each $C_\xi\subseteq B_\xi$ has value zero: multiplying its idempotent
by $h_{B_\xi}$ leaves it unchanged. By \eqref{eq:smallidempotent},
$\bigcup_{\xi<\mu}C_\xi=\bigcup_{\xi<\mu}B_\xi$ also has value zero.
Its complement $\bigcap_{\xi<\mu}A_\xi$ is therefore in $U$.
Thus $U$ is a nonprincipal $\kappa$-complete ultrafilter on $\kappa$,
which is the usual measure characterization of measurability.
\end{proof}

This criterion detects measurability only because the mask algebra and its
small-union condition explicitly retain a $\kappa$-indexed Boolean algebra.
It is not a characterization of measurability in the pure ordered-field
language of $\No$.

\section{Omnific integers and a purely infinite measure code}\label{sec:omnific}

\subsection{The ring and its short-support part}
The omnific integers form the discretely ordered subring
\begin{equation}\label{eq:Oz}
 \Oz=\left\{\sum_a r_a\omega^a:
      \supp(x)\subseteq\No^{\ge0},\quad r_0\in\Z\right\}.
\end{equation}
A missing zero exponent has coefficient zero. Nonzero coefficients at
positive exponents are arbitrary reals, not necessarily ordinary integers.
This is Conway's normal-form description \cite{Conway,Gonshor}.
We use it rather than an informal criterion such as ``having no fractional
part,'' which can hide a convention at exponent zero.
Put
\[
 \Ik=\Oz\cap\Sk.
\]

\begin{theorem}[Omnific equalizer, intersection, and fractions]\label{thm:Oz}
The maps $J$ and $H_j$ preserve and reflect membership in $\Oz$, and
\begin{align}
 \{z\in\Oz:J(z)=H_j(z)\}&=\Ik,\label{eq:Ozeq}\\
 J(\Oz)\cap H_j(\Oz)&=J(\Ik)=H_j(\Ik),\label{eq:Ozinter}\\
 \Frac(\Ik)&=\Sk,\label{eq:Ozfrac}\\
 \Frac\bigl(J(\Oz)\cap H_j(\Oz)\bigr)&=\Common.\label{eq:commonfrac}
\end{align}
Also $\ct(J(z))=\ct(H_j(z))=\ct(z)$ for every $z\in\Oz$.
\end{theorem}
\begin{proof}
For $H_j$ this preservation and reflection is part of
\leref{lem:reindex}. For $J$, transfer the canonical normal form and the
condition in \eqref{eq:Oz}. Nonnegativity of all exponents and the integer
condition on the constant coefficient are absolute; equivalently, apply
elementarity to the normal-form description and then
\leref{lem:absolute}. A coefficient at exponent zero is fixed, and a new
missing-image exponent cannot equal zero because $0\in\Gam$.
The same reasoning gives the constant-term identity for both maps.

The equalizer formula is \thref{thm:defect} restricted to $\Oz$.
If $y\in J(\Oz)\cap H_j(\Oz)$, write $y=J(z)$ with $z\in\Oz$.
The field intersection theorem forces $z\in\Sk$, so $z\in\Ik$.
Conversely, both maps have the same value on any $z\in\Ik$.
This proves \eqref{eq:Ozinter}.

To prove \eqref{eq:Ozfrac}, take nonzero $x\in\Sk$. Its support is a
set, so choose $d>0$ such that $d>-a$ for every $a\in\supp(x)$.
The monomial $q=\omega^d$ and the shifted series $p=\omega^d x$ have
only strictly positive exponents. Their constant coefficients vanish,
and both are in $\Oz$. Their supports have cardinality less than
$\kappa$, so $p,q\in\Ik$ and $x=p/q$. The reverse inclusion follows
because $\Sk$ is a field containing $\Ik$. Applying the field embedding
$J$ gives \eqref{eq:commonfrac}.
\end{proof}

\subsection{A defect with no finite residue}
For $A\subseteq\kappa$, define
\begin{equation}\label{eq:zA}
 z_A=\omega^\kappa h_A
     =\sum_{\alpha\in A}\omega^{\kappa-\alpha}.
\end{equation}
Every exponent is strictly positive, since $\alpha<\kappa$.
Thus $z_A$ is a purely infinite omnific integer when $A$ is nonempty,
and is zero otherwise. Here and below $\kappa-\alpha$ is the difference
in the surreal ordered field.

\begin{theorem}[Omnific recovery of the measure]\label{thm:Ozmeasure}
For $A\subseteq\kappa$,
\begin{align}
 D_j(z_A)&=\omega^{j(\kappa)}D_j(h_A),\label{eq:zdefect}\\
 \coef{j(\kappa)-\kappa}{D_j(z_A)}
      &=\mathbf 1_{\{A\in U_j\}},\label{eq:zmeasure}\\
 \ct(D_j(z_A))&=0.\label{eq:zct}
\end{align}
In particular, for $z_\kappa=\omega^\kappa h_\kappa$,
\[
 D_j(z_\kappa)=
   \sum_{\kappa\le\beta<j(\kappa)}\omega^{j(\kappa)-\beta}>0,
 \qquad
 \LT(D_j(z_\kappa))=\omega^{j(\kappa)-\kappa}.
\]
Every nonzero term of every $D_j(z_A)$ is purely infinite.
\end{theorem}
\begin{proof}
The two embeddings agree on $\omega^\kappa$, with common value
$\omega^{j(\kappa)}$, and both are multiplicative. This proves
\eqref{eq:zdefect}. Multiplication by that monomial shifts exponents by
$j(\kappa)$, so \eqref{eq:zmeasure} follows from
\eqref{eq:measurecoef}. Moreover $j(A)\cap\kappa=A$, so the defect of
$h_A$ has its possible indices only in
$j(A)\setminus A\subseteq[\kappa,j(\kappa))$.
After shifting, all exponents $j(\kappa)-\beta$ are strictly positive,
which proves \eqref{eq:zct} and the last assertion. Taking $A=\kappa$
gives the explicit nonempty series and its leading term.
\end{proof}

The constant-term map therefore loses this large-cardinal information,
even though one specified infinite-scale coefficient recovers it.
This is compatible with the repository's emphasis on constant-term
factorizations in omnific arithmetic \cite{RepoQuotients}; it does not
supply a new factorization theorem for set-sized quotient rings.

\section{Surcomplex numbers and the Gaussian omnific lattice}\label{sec:complex}

Let $\SC=\No(i)$ with a fixed named element $i$ satisfying $i^2=-1$,
and write conjugation as $\overline{x+iy}=x-iy$.
A surcomplex normal form has complex coefficients and a reverse
well-ordered set of surreal exponents. Its support is the union of the
supports of its real and imaginary coordinates: a complex coefficient
vanishes exactly when both real coordinates vanish.
Define
\[
 \Jc(x+iy)=J(x)+iJ(y),\qquad
 \Hc(x+iy)=H_j(x)+iH_j(y).
\]
These definitions fix the chosen orientation $i$; they are not a statement
about arbitrary automorphisms of the pure surcomplex field.

\begin{theorem}[Surcomplex transfer]\label{thm:complex}
The maps $\Jc,\Hc:\SC\to\SC$ are field embeddings fixing $\C$ and
commuting with conjugation. The second is strongly $\C$-linear.
They are elementary in the pure field language, and
\begin{align}
 \operatorname{Eq}(\Jc,\Hc)&=\Sk(i),\label{eq:complexeq}\\
 \Jc(\SC)\cap\Hc(\SC)&=\Common(i).\label{eq:complexinter}
\end{align}
Both fields on the right are algebraically closed. The image fields on
the left of \eqref{eq:complexinter} are incomparable. Their intersection
with the Gaussian omnific lattice satisfies
\begin{equation}\label{eq:Gaussian}
 \Jc(\Oz[i])\cap\Hc(\Oz[i])=J(\Ik)[i].
\end{equation}
\end{theorem}
\begin{proof}
The quadratic-extension multiplication law proves the field-homomorphism
claims coordinatewise. Uniqueness of real and imaginary parts gives
injectivity and commutation with conjugation. The strong summation claim
is the real statement on both coordinates. Since $\No$ is real closed,
$\SC$ and both image fields are algebraically closed of characteristic
zero. Quantifier elimination for $\ACF$ gives elementarity, formula by
formula, as before.

Equality of $\Jc(x+iy)$ and $\Hc(x+iy)$ is equivalent to both
$J(x)=H_j(x)$ and $J(y)=H_j(y)$. This is precisely $x,y\in\Sk$.
Equivalently, the union of the two supports has size less than $\kappa$:
a finite union of sets of size less than an infinite cardinal still has
size less than that cardinal. This proves \eqref{eq:complexeq}.
The same coordinate comparison and \thref{thm:intersection} prove
\eqref{eq:complexinter}. The fields $\Sk$ and $\Common$ are real closed,
so adjoining $i$ makes them algebraically closed.

The real witnesses $h_\kappa$ and $h_{j(\kappa)}$ still witness
incomparability. If a real element is the image of $x+iy$, its imaginary
coordinate is zero, forcing $y=0$; allowing a complex preimage cannot
repair the failure of a real preimage. Finally apply
\thref{thm:Oz} separately to the two integer coordinates to obtain
\eqref{eq:Gaussian}.
\end{proof}

The exact defect formula also holds directly with complex coefficients:
for $z=\sum_{\alpha<\theta}c_\alpha\omega^{a_\alpha}$, set
$c^*=j(c)$ after identifying $\C$ with $\R^2$. Then
\[
 \Jc(z)-\Hc(z)=
   \sum_{\beta\in P_\theta}c^*_\beta\omega^{a^*_\beta}.
\]
Nonzero complex coefficients at distinct exponents cannot cancel. Thus
long complex support, not a sign or order on complex numbers, is what
forces disagreement. The Boolean measure masks from
Section~\ref{sec:measure} already lie on the real axis and require no
complex choice beyond the named inclusion.

\begin{corollary}[Conjugation and modulus]\label{cor:modulus}
If $|z|=(z\overline z)^{1/2}\in\No^{\ge0}$, then
\[
 |\Jc(z)|=J(|z|),\qquad |\Hc(z)|=H_j(|z|).
\]
Every surcomplex number in either image outside $\Common(i)$ is
transcendental over $\Common(i)$.
\end{corollary}
\begin{proof}
Both maps preserve conjugation and the unique nonnegative square root.
The transcendence assertion follows because $\Common(i)$ is algebraically
closed.
\end{proof}

\section{Defect algebra and a logarithmic homomorphism}\label{sec:log}

\subsection{What all old-image coefficients fail to see}
\begin{proposition}[Coefficient invisibility]\label{prop:invisible}
For every $a,x\in\No$,
\begin{equation}\label{eq:oldcoeff}
 \coef{J(a)}{J(x)}=\coef{a}{x}=\coef{J(a)}{H_j(x)}.
\end{equation}
In particular, $J$ and $H_j$ preserve the coefficient at exponent zero.
They preserve the finite-surreal valuation ring and its infinitesimal
ideal, and induce the identity on its residue field $\R$.
\end{proposition}
\begin{proof}
If $a$ is in the support of $x$, its unique original index gives the
coefficient in \eqref{eq:oldindices}. A missing-image index cannot supply
exponent $J(a)$ by \thref{thm:defect}. If $a$ is not in the support, the
same range-reflection argument used in that theorem excludes $J(a)$ from
the support of $J(x)$. The identity for $H_j$ is its definition.
Setting $a=0$ gives the constant coefficient. The valuation identities
in Corollary~\ref{cor:leading}, together with $J(0)=0$ and order
preservation, give preservation of the valuation ring and its maximal
ideal. On a finite surreal, its residue is its coefficient at zero.
\end{proof}

Thus \eqref{eq:measurecoef} detects genuinely new exponent positions. No
coefficient extraction at an exponent in $J(\No)$ can see the defect.
Both embeddings nevertheless have the same value-group action and the
same residue action.

\begin{proposition}[Twisted product law]\label{prop:twisted}
The additive map $D_j$ satisfies
\begin{equation}\label{eq:twisted}
 D_j(xy)=J(x)D_j(y)+D_j(x)H_j(y).
\end{equation}
For $c\in\Sk$ one has $D_j(cx)=J(c)D_j(x)$.
\end{proposition}
\begin{proof}
Expand the right side as
$J(x)J(y)-J(x)H_j(y)+J(x)H_j(y)-H_j(x)H_j(y)$.
The middle terms cancel and multiplicativity gives the left side.
The scalar identity follows because $D_j(c)=0$.
\end{proof}

Equation~\eqref{eq:twisted} is a twisted derivation identity, not an
ordinary Leibniz rule for a derivation on $\No$. No claim is made that
$D_j$ is a Berarducci--Mantova derivation or a strongly additive derivation.

\subsection{Passing from a multiplicative discrepancy to an additive one}
For $x\ne0$, let
\begin{equation}\label{eq:eta}
 \eta_j(x)=\frac{J(x)}{H_j(x)}.
\end{equation}
By Corollary~\ref{cor:leading}, this ratio lies in $1+\infideal$.
For $u\in\infideal$, the formal series
\begin{equation}\label{eq:logseries}
 \log(1+u)=\sum_{n\ge1}\frac{(-1)^{n+1}}{n}u^n
\end{equation}
is strongly summable by the Hahn--Neumann lemma. The usual formal
power-series identities are legitimate after such infinitesimal
substitution; in particular, this logarithm takes products in
$1+\infideal$ to sums \cite{Gonshor,vdDE}. Only this near-one logarithm
is needed, not a preservation theorem for the global surreal exponential.

\begin{theorem}[Logarithmic defect]\label{thm:logdefect}
The map
\[
 L_j:\No^\times\longrightarrow(\infideal,+),\qquad
 L_j(x)=\log\eta_j(x)
\]
is a group homomorphism with kernel $\Sk^\times$. If $x$ has normal form
\eqref{eq:NF} of length $\theta\ge\kappa$, then
\begin{equation}\label{eq:LTlog}
 \LT(L_j(x))=
   \frac{r^*_{\kappa}}{r_0}\,
       \omega^{a^*_{\kappa}-J(a_0)}.
\end{equation}
Thus a nonzero additive infinitesimal invariant records every departure
from the equalizer, after multiplicative normalization.
\end{theorem}
\begin{proof}
Because both embeddings are multiplicative,
$\eta_j(xy)=\eta_j(x)\eta_j(y)$. The logarithm identity gives
$L_j(xy)=L_j(x)+L_j(y)$.
If $\eta_j(x)=1+u$ with $u\ne0$, then the first term of
\eqref{eq:logseries} is $u$, while all higher powers have strictly larger
valuation. Hence the logarithm is nonzero and has leading term $\LT(u)$.
It follows that $L_j(x)=0$ precisely when $J(x)=H_j(x)$, proving the
kernel assertion.

Finally $u=D_j(x)/H_j(x)$. Its leading term is the quotient of the
leading term in \eqref{eq:firstterm} by the common leading term from
Corollary~\ref{cor:leading}, namely \eqref{eq:LTlog}. Its exponent is
negative since $a^*_{\kappa}<a^*_0=J(a_0)$, so it is infinitesimal as
required.
\end{proof}

For the basic witness,
\[
 \LT(L_j(h_\kappa))=\omega^{-\kappa}.
\]
The homomorphism formulation avoids introducing a quotient whose
individual equivalence classes might be proper classes. We use its
kernel as a class predicate; no extra axiom for forming hyperclasses or
class-sized quotient objects is needed.

\section{An exact exponential compatibility locus}\label{sec:exponential}
This section adds the canonical Gonshor exponential. It is independent of
the proof of the main defect and image-intersection theorems, but the same
support calculation yields a second exact boundary.

Write every surreal uniquely as
\[
 x=p(x)+r(x)+\eps(x),
\]
where $p(x)$ has only positive exponents, $r(x)\in\R$, and
$\eps(x)\in\infideal$. Thus $p(x)$ is its purely infinite part.
We use the classical exponential facts that $\exp(p)$ is a Conway monomial
for purely infinite $p$, that the exponential on infinitesimals is its
strong power series, and that the inverse logarithm of a monomial is
purely infinite \cite[Chapter~10]{Gonshor}\cite[Section~5]{vdDE}.

\begin{lemma}[Exponential transfer]\label{lem:exptransfer}
The elementary-induced embedding satisfies
\[
 J(\exp x)=\exp J(x),\qquad
 J(\log y)=\log J(y)\quad(y>0).
\]
\end{lemma}
\begin{proof}
Gonshor's exponential is defined by a uniform recursive cut on the
canonical options of its argument, using exponentials of simpler options,
finite Taylor polynomials, field operations, and their inequalities.
The cut absoluteness argument from \leref{lem:absolute} applies to this
recursion as well. Therefore the exponential computed in $M$ on an old
argument equals the one in $V$. Elementary transfer of the defining
recursion gives the first identity. The second follows from the uniqueness
of the inverse of the strictly increasing bijection $\exp$.
\end{proof}

\begin{theorem}[Exact exponential defect]\label{thm:expdefect}
For every $x\in\No$,
\begin{equation}\label{eq:expdefect}
 \frac{H_j(\exp x)}{\exp(H_j(x))}
      =\exp\bigl(D_j(p(x))\bigr).
\end{equation}
In particular,
\begin{equation}\label{eq:explocus}
 H_j(\exp x)=\exp(H_j(x))
    \quad\Longleftrightarrow\quad
 \abs{\supp(p(x))}<\kappa.
\end{equation}
The companion commutes with the exponential on every finite surreal,
including finite surreals with arbitrarily long infinitesimal support.
For a positive $y=r\omega^a(1+u)$, with $r>0$ real and $u\in\infideal$,
\begin{equation}\label{eq:logcompat}
 \log(H_j(y))-H_j(\log y)=D_j(\log(\omega^a)).
\end{equation}
\end{theorem}
\begin{proof}
The companion preserves the purely infinite, real, and infinitesimal
parts separately, because it preserves the sign of exponents. Since
$\exp(p(x))$ is a single monomial,
\[
 H_j(\exp(p(x)))=J(\exp(p(x)))=\exp(J(p(x))).
\]
Strong linearity and infinitesimal power-series evaluation give
$H_j(\exp(\eps(x)))=\exp(H_j(\eps(x)))$, while real exponential
values are fixed. Applying $H_j$ to
$\exp x=\exp(p(x))e^{r(x)}\exp(\eps(x))$ and dividing by
$\exp(H_j(x))$ proves \eqref{eq:expdefect}. Injectivity of the
exponential and \thref{thm:defect} give \eqref{eq:explocus}.
For finite $x$, its purely infinite part is zero.

Finally use
$\log y=\log r+\log(\omega^a)+\log(1+u)$.
The near-one logarithm commutes with $H_j$ by its strong power-series
formula. Also
$\log(H_j(\omega^a))=\log(J(\omega^a))=J(\log(\omega^a))$.
Subtracting the two decompositions gives \eqref{eq:logcompat}.
\end{proof}

\begin{corollary}[Exponential closure without logarithmic closure]\label{cor:expclosure}
The fields $\Sk$ and $\Common$ are closed under the Gonshor exponential
but not under logarithms of positive elements. With this induced
exponential, neither is an elementary substructure of $(\No,\exp)$,
although both are elementary ordered subfields of $\No$.
\end{corollary}
\begin{proof}
If $x\in\Sk$, then $p(x)\in\Sk$. By the preceding theorem and
\leref{lem:exptransfer},
\[
 H_j(\exp x)=\exp(H_j(x))=\exp(J(x))=J(\exp x).
\]
The equalizer theorem yields $\exp x\in\Sk$. The same closure for
$\Common=J(\Sk)$ follows from exponential transfer.

The purely infinite element $z_\kappa$ of Section~\ref{sec:omnific} has
$\kappa$ outer terms, so it is not in $\Sk$. But $\exp(z_\kappa)$ is
a positive monomial and hence belongs to $\Sk$. Its unique logarithm
in $\No$ is $z_\kappa$, which is absent. Similarly
$J(\exp(z_\kappa))$ lies in $\Common$, whereas its logarithm
$J(z_\kappa)$ does not, by injectivity of $J$ and the description of
$\Common$. Thus the sentence asserting that every positive element has
an exponential preimage holds in $(\No,\exp)$ and fails in both smaller
exponential structures. Their ordered-field elementarity follows from
real closedness.
\end{proof}

\begin{corollary}[One image is exponentially closed and the other is not]
\label{cor:Hnotexp}
The image $J(\No)$ is closed under both exponential and logarithm.
The image $H_j(\No)$ is not closed under exponential.
\end{corollary}
\begin{proof}
The first assertion follows immediately from \leref{lem:exptransfer}.
For the second, use the classical order-preserving bijection
$g:\No^{>0}\to\No$ specified by
\[
 \exp(\omega^b)=\omega^{\omega^{g(b)}}\quad(b>0).
\]
The purely infinite exponential formula says that
$\exp(\sum r_\alpha\omega^{b_\alpha})
=\omega^{\sum r_\alpha\omega^{g(b_\alpha)}}$ when all $b_\alpha>0$;
see \cite[Theorem~10.13]{Gonshor} and \cite[Lemma~5.2, proof]{vdDE}.
Define
\[
 p_\kappa=\sum_{\alpha<\kappa}\omega^{g^{-1}(-\alpha)}.
\]
The exponents are strictly decreasing and positive, so this is a purely
infinite normal form and $\exp(p_\kappa)=\omega^{h_\kappa}$.
The definition of $g$, together with transfer of exponential and the
omega-map, shows that $J(g(b))=g(J(b))$ and hence that $J$ commutes with
$g^{-1}$. Since $-\alpha$ is fixed for $\alpha<\kappa$, all the displayed
exponents are fixed. Therefore $H_j(p_\kappa)=p_\kappa$, so
$p_\kappa\in H_j(\No)$.

But $h_\kappa\notin J(\No)=\Gam$ by \thref{thm:witness}.
The monomial $\omega^{h_\kappa}$ consequently has its exponent outside
$\Gam$ and is not in $H_j(\No)$ by \eqref{eq:Himage}.
This proves the failure of exponential closure.
\end{proof}

\begin{theorem}[Two canonical closures of the common field]\label{thm:twoclosures}
Among class subfields of $\No$ containing $\Common$, the least one closed
under logarithms of positive elements is $J(\No)$, and the least one
closed under external set-indexed strong sums is $H_j(\No)$.
More explicitly, every element of $J(\No)$ has the form
\begin{equation}\label{eq:singlelog}
 \frac{\log m_1}{m_2},
 \qquad m_1,m_2\in\{\omega^{J(a)}:a\in\No\}\subseteq\Common,
\end{equation}
and every element of $H_j(\No)$ is a strong sum of elements of
$\Common$. These two closures are incomparable and intersect in exactly
$\Common$.
In contrast, both the logarithmic field closure and the strong-sum closure
of the domain equalizer $\Sk$ are all of $\No$.
\end{theorem}
\begin{proof}
The field $J(\No)$ contains $\Common$ and is logarithmically closed by
\leref{lem:exptransfer}. Take $J(x)\in J(\No)$.
Choose $d>0$ so that $z=\omega^d x$ is purely infinite, exactly as in
\thref{thm:Oz}. Then $\exp z$ is a monomial. Both $\exp z$ and
$\omega^d$ have one-term support and hence lie in $\Sk$.
Set $m_1=J(\exp z)$ and $m_2=J(\omega^d)$; they are monomials in
$\Common$. Exponential transfer gives
\[
 J(x)=\frac{J(z)}{J(\omega^d)}=\frac{\log m_1}{m_2}.
\]
Any logarithmically closed class subfield containing $\Common$ must
therefore contain every $J(x)$, proving minimality and
\eqref{eq:singlelog}.

The field $H_j(\No)$ is strongly closed by \thref{thm:escape}. Every
normal-form term $r\omega^{J(a)}$ in one of its members already lies in
$\Common$, so strong closure of any class containing $\Common$ forces
that member to be included. This proves the other minimality assertion.
Incomparability and the intersection were proved earlier.

For $\Sk$, the identical shift argument without $J$ expresses every
$x\in\No$ as a quotient of the logarithm of a monomial in $\Sk$ by
another monomial in $\Sk$. Also every surreal is a strong sum of real
multiples of monomials, all of which lie in $\Sk$. This proves both
final assertions. The minimality statements are explicit universal
properties; they do not require forming an intersection of a collection
of all classes.
\end{proof}

The two image fields are therefore canonically isomorphic as ordered
valued fields over $\Common$, but their actual placement in $\No$ has
different compatibility with the canonical exponential. Neither their
shared residue field and value group nor their common monomials detect
this difference. The support criterion \eqref{eq:explocus} is different
from the whole-support equalizer criterion: the long-support finite
number $h_\kappa$ lies in the former locus and not the latter.

\section{Supercompactness as a seed-coefficient system}\label{sec:supercompact}

The first support threshold depends only on the critical point. Thus it
cannot by itself distinguish a measurable embedding from a stronger
embedding with the same critical point. A richer family of coefficient
probes does retain the classical seed information used for
supercompactness.

Assume $\lambda\ge\kappa$, $j(\kappa)>\lambda$, and $M$ is closed under
$\lambda$-sequences from $V$. Let
\[
 I=P_\kappa(\lambda)=\{a\subseteq\lambda:\abs a<\kappa\},\qquad
 s=\imset j\lambda.
\]
Closure gives $s\in M$, and in $M$ it has size at most $\lambda<j(\kappa)$,
so $s\in j(I)$. These are the standard hypotheses and seed for a
$\lambda$-supercompactness measure \cite{Neeman}.
Choose a bijection $e:\theta\to I$. Let $\beta_s<j(\theta)$ be the
unique index satisfying
\begin{equation}\label{eq:seedindex}
 j(e)(\beta_s)=s.
\end{equation}
For $A\subseteq I$, define the mask
\[
 m_A=\sum_{\substack{\alpha<\theta\\e(\alpha)\in A}}\omega^{-\alpha}.
\]
The construction depends on the enumeration $e$, which is explicitly part
of the data.

\begin{theorem}[Seed extraction]\label{thm:supercompact}
The index $\beta_s$ is a missing-image index:
$\beta_s\notin\imset j\theta$. Moreover
\begin{equation}\label{eq:seedcoef}
 \coef{-\beta_s}{D_j(m_A)}
     =\mathbf 1_{\{s\in j(A)\}}\qquad(A\subseteq I).
\end{equation}
Consequently the seed coefficient recovers the fine normal
$\kappa$-complete ultrafilter
\[
 U_s=\{A\subseteq P_\kappa(\lambda):s\in j(A)\}.
\]
\end{theorem}
\begin{proof}
If $\beta_s=j(\alpha)$, then $s=j(e(\alpha))=j(a)$ for some
$a\in P_\kappa(\lambda)$. Since $\abs a<\kappa$,
\leref{lem:index} gives $j(a)=\imset ja$. But
$s=\imset j\lambda$, so injectivity of $j$ implies $a=\lambda$, a
contradiction. Thus $\beta_s$ is missing from the pointwise index image.

The transformed mask has coefficient 1 at $-\beta_s$ exactly when
$j(e)(\beta_s)=s$ belongs to $j(A)$. The companion has its exponents
only at positions $-j(\alpha)$, so it contributes zero there. This gives
\eqref{eq:seedcoef}.

Ultrafilter and $\kappa$-completeness follow by the same Boolean and
short-index argument as in \thref{thm:measure}, with $s$ replacing
$\kappa$. For each $\xi<\lambda$, the seed contains $j(\xi)$, so
\[
 \{a\in P_\kappa(\lambda):\xi\in a\}\in U_s;
\]
this is fineness. If $A\in U_s$ and $f:A\to\lambda$ satisfies
$f(a)\in a$, then $j(f)(s)\in s=\imset j\lambda$.
Write $j(f)(s)=j(\xi)$ for some $\xi<\lambda$. The fiber where
$f(a)=\xi$ then belongs to $U_s$, proving normality.
\end{proof}

Fine normal $\kappa$-complete ultrafilters on $P_\kappa(\lambda)$ are the
classical ultrafilter characterization of supercompactness, due to Solovay;
see \cite{Neeman}. We have not discovered a new large-cardinal equivalence.
The contribution is that the seed can be read as a coefficient of the
same defect construction, and that its index is provably outside the
termwise image. An intrinsic, enumeration-free coefficient geometry for
these seed systems is a further research question, not a result claimed
here.

\section{What large cardinals change, and what they do not}\label{sec:interpretation}

\subsection{No change to finite algebraic truth}
Both $J$ and $H_j$ are elementary ordered-field embeddings. Their complex
lifts are elementary embeddings of algebraically closed fields. Consequently
they preserve every first-order formula in those respective field
languages. The distinction made in this article is in additional structure:
normal forms, external supports, indexed strong sums, and their coefficient
probes. It would be wrong to describe \eqref{eq:Dh} as a failure of an
ordinary field axiom.

Likewise, a measurable cardinal does not alter the formula for adding two
specified old surreal sign sequences. It supplies a nontrivial elementary
embedding of the surrounding set universe, and that map moves certain
\emph{entire set-indexed presentations}. The finite identities still
transfer; the pointwise index set does not exhaust the transferred one.

\subsection{Why there is no Kunen contradiction}
The set-theoretic map is $j:V\to M$, not an elementary self-embedding
$V\to V$. After forming surreal numbers in $M$, we compose with the
identity inclusion into the ambient $\No^V$. That inclusion is elementary
in the \emph{ordered-field reduct}, by real closedness. It is not being
claimed elementary in a structure that recovers all of the surrounding
set theory from birthdays or other codes.

This matters when comparing with the repository's birthday-enriched
structures \cite{RepoBirthday}. An order-field embedding can preserve
specific operations and a birthday identity on its image without making
the full ambient birthday expansion an elementary extension. The extra
quantifiers of that expansion can range over new sign codes. Accordingly,
our construction is not an exception to Kunen-type obstructions for a
set-theoretically faithful elementary self-embedding.

\subsection{A hierarchy of information}
The first failure cardinal of positive strong summation is $\kappa$.
The first missing coefficient on the masks $h_A$ recovers the normal
measure $U_j$. At a supercompact seed, the wider mask family $m_A$
recovers the associated fine normal measure. These are three different
levels of data:
\[
 \text{critical point}
 \quad\longleftarrow\quad
 \text{normal-measure coefficient functional}
 \quad\longleftarrow\quad
 \text{seed-coefficient systems}.
\]
The arrows mean forgetting information, not an implication that every
normal measure extends to every seed system. In particular, two distinct
normal measures can have the same critical cardinal and the same equalizer
field $\Sk$, while their defect functionals differ by
Corollary~\ref{cor:distinctU}.

\subsection{Choice and the scope of set formation}
The proof does use Choice, already present in $\ZFC$, to enumerate an
arbitrary set in \leref{lem:index} and to use the conventional measure
characterizations. It uses Replacement with the embedding parameter to
form pointwise images and inverse images of \emph{sets}. It does not
choose a basis for the entire proper class $\No$, does not appeal to a
class version of Zorn's lemma, and does not form the collection of all
class endomorphisms as a set or class of its own. A definable normal
ultrapower provides all the specific maps needed.

\section{Position relative to the literature and the supplied repository}
\label{sec:position}

\subsection{Classical ingredients and proposed contributions}
The following separation is part of the claim, not merely an editorial
qualification. Conway normal forms, the real closedness of $\No$, the
algebraic closedness of $\No(i)$, Hahn reindexing, and Gonshor's
exponential are established background. Normal measures derived from
embeddings, the ultrapower construction, and the fine-normal-measure
characterization of supercompactness are also established background.
The article uses these results to obtain an explicit comparison inside
one ambient surreal field.

\begin{longtable}{@{}p{0.26\textwidth}p{0.36\textwidth}p{0.32\textwidth}@{}}
\toprule
\textbf{Item} & \textbf{Status in this article} & \textbf{Relevant precedent}\\
\midrule\endhead
Conway normal forms and strong summation & Classical input; conventions and
absoluteness requirements made explicit. & Conway, Gonshor, van den
Dries--Ehrlich, Ehrlich \cite{Conway,Gonshor,vdDE,EhrlichNames}.\\[5pt]
Monomial reindexing & Classical Hahn construction; elementary proof supplied.
& General Hahn-field theory and the strong maps studied in \cite{KKS}.\\[5pt]
Critical-index defect & Candidate original exact formula, including its first
term and exclusion of image exponents. & Standard distinction
$j(I)\ne\imset jI$ supplies the set-theoretic mechanism.\\[5pt]
Equalizer and intersection & Candidate original package: exact short-support
equalizer, common-image formula, incomparability, and coefficient projection.
& Real-closed-field facts are classical; the present pair is constructed
from the embedding.\\[5pt]
Strong-sum escape & Candidate original intrinsic closure threshold for
$J(\No)$, not only a noncommutation example. & Strong summability itself is
classical.\\[5pt]
Measure coefficient & Candidate original arithmetic realization. The derived
ultrafilter and its properties are not new. & Normal ultrapowers and seed
measures \cite{Neeman}.\\[5pt]
Mask criterion & Classical ultrafilter criterion transported to an explicitly
specified Hadamard algebra. & Direct-product homomorphisms and ultrafilters
\cite{Bergman}.\\[5pt]
Exponential boundary & Candidate original compatibility locus for this
companion and its two-closure description. & Gonshor's exponential
normal-form formula \cite{Gonshor,vdDE}.\\[5pt]
Omnific and surcomplex transfer & Consequences proved here from the real
construction; no separate large-cardinal assumption. & Canonical integer
parts and quadratic extensions are classical.\\
\bottomrule
\end{longtable}

The phrase ``candidate original'' means that this formulation and proof
package was not located in the targeted material inspected for this
article. It does not assert that no specialist has previously observed
some or all of the consequences. In particular, elementary embeddings
and Hahn series are mature subjects; the mechanism is simple once the
correct two images are compared.

\subsection{The repository snapshot actually inspected}
The repository was read at commit
\begin{center}\small
\texttt{0865f043aec113c14c69ef45006bbc7546a4e75a}.
\end{center}
The targeted inspection comprised the root \texttt{README.md}, the
\texttt{docs/README.md} report index, the READMEs for
\emph{Surreal Fields Across Set-Theoretic Universes} and
\emph{Birthday Cutoffs and Hereditary Sets}, and the former article's
LaTeX preamble \cite{Repo,RepoUniverses,RepoBirthday}.
The style here follows that report's restrained theorem-centered layout.
This was not a line-by-line audit of the whole repository or its formalization
ledger.

The cross-universe report describes omitted cuts, fresh sign sequences,
external saturation, first new birthdays, and Prikry examples. The
birthday report describes hereditary-set interpretations and the recovery
of set-theoretic information from birthday-enriched structures. Its README
explicitly attributes the announced full-class bi-interpretation to
Chen--Hamkins--Yang. Those earlier themes are relevant context but are not
used as unproved premises of our theorems.

The present article takes a different route. It does not add birthday
coding to recover the entire set universe, and it does not change the
universe by forcing. It compares the pointwise elementary action and its
strong monomial extension inside a fixed universe. The output is an
arithmetic defect whose support and coefficients can be calculated exactly.
The constant-term discussion of omnific integers is also kept separate
from the repository's proposed set-sized quotient theorems
\cite{RepoQuotients}.

The root README reports existing formal Hahn infrastructure, including
summable families and exponent-embedding maps. That makes a substantial
part of the algebraic proof a plausible formalization target. It does
\emph{not} mean the large-cardinal construction, the absoluteness lemma,
or any theorem specific to this article is already checked in Lean.
The repository itself warns that research reports, finite tests, and
formal proof coverage are different kinds of evidence.

\subsection{Relation to current automorphism work}
Kaplan--Krapp--Serra \cite{KKS} provide a contemporary framework for
strongly additive maps, monomial actions, valuation structure, and
additional surreal operations. We use their definitions and general Hahn
context, not a particular non-strong automorphism construction from that
paper. Our maps are nonsurjective embeddings, and their valuations are
transported by $J$ rather than fixed pointwise. Results that assume
bijectivity or pointwise preservation of the value group must not be
applied to this pair without checking those hypotheses.

No claim is made to settle one of that paper's numbered questions.
Rather, our construction answers a new precisely stated family of
structural questions under an explicit large-cardinal hypothesis. A
specialist comparison with the unpublished details of related projects
would be appropriate before submitting a priority claim.

\section{Further research questions and concrete next steps}\label{sec:questions}

The following questions are proposed by this article. They are not
presented as an established list of famous open problems. Each isolates
something not proved above and indicates what a meaningful answer would
require.

\begin{question}[Optimal consistency strength]\label{q:strength}
Suppose two ordered-field embeddings $P,Q:\No\to\No$ fix $\R$, agree
on all ordinals and monomials, $Q$ is strongly linear, and their equalizer
and image intersection satisfy the formulas of
\thref{thm:intersection} for an uncountable $\kappa$. Suppose also that
$P$ fixes every ordinal below $\kappa$, moves $\kappa$, and commutes with
the omega-map. Must $\kappa$ be measurable, or can such a pair exist
without an ambient elementary embedding of set theory?
\end{question}
The obstruction to a converse is precise: the first defect coefficient
is not automatically a character for the Hadamard product merely because
$P$ is multiplicative for ordinary surreal multiplication.
\thref{thm:measurablecriterion} solves the converse only when the
coefficient-algebra structure and its small-union law are available.

\begin{question}[Classification of strong companions]
For an arbitrary real-fixing omega-map-compatible ordered-field embedding
$P$, classify the equalizer of $P$ with the Hahn reindexing map $T_P$.
Which equalizers are cardinal support bounds, and which depend on the
shape, cofinality, or iterated structure of the exponent support?
\end{question}
The construction $T_P$ is always available by \leref{lem:reindex}.
What is special to the present large-cardinal setting is the exact
missing-index formula and the absence of cancellation in it.

\begin{question}[Recovering an embedding from fewer coefficient probes]
What is the least family of masks and coefficient positions that determines
$j$ on a specified hereditary-size universe $H_\lambda$? Can normal-form
probes of bounded complexity recover more than the derived normal measure
at the critical point?
\end{question}
Using all subsets of all ordinals with all coefficients would retain
far more information than a single normal measure. The meaningful task
is to specify a restricted family and prove a sharp reconstruction or
nonreconstruction theorem.

\begin{question}[Measure order from defect data]
Can the Mitchell ordering, or another comparison relation between normal
measures on $\kappa$, be recovered from a natural algebraic relation between
their coefficient functionals $\chi_j$, without naming the ambient
ultrapower models as extra predicates?
\end{question}
Corollary~\ref{cor:distinctU} gives injective detection of individual
measures. It gives no criterion for their relative position in an
ultrapower hierarchy.

\begin{question}[Intrinsic supercompact seed systems]
Can the seed construction of \thref{thm:supercompact} be formulated
without choosing an enumeration of $P_\kappa(\lambda)$? Is there a natural
coordinate-free compatibility law for the resulting systems as $\lambda$
varies that detects the appropriate closure of $M$?
\end{question}
Any proposal must preserve fineness and normality and must distinguish
an actual seed in $j(P_\kappa(\lambda))$ from an arbitrary missing-image
index. Merely having many extra coefficients does not imply
supercompactness.

\begin{question}[Joint algebraic dependence of the images]
Are $J(\No)$ and $H_j(\No)$ linearly disjoint over $\Common$? More
modestly, for a finite tuple $x$, can one decide the transcendence degree
of $H_j(\No)(J(x))$ over $H_j(\No)$ from the missing-image supports?
\end{question}
Every single element outside the intersection is transcendental over
the opposite real-closed image field. That does not decide joint algebraic
independence of a tuple, and the exact intersection theorem is not a
substitute for a disjointness proof.

\begin{question}[Alternating the two closures]
Starting with $\Common$, what field is obtained by alternating logarithmic
field closure and external strong-sum closure? Do finitely many alternations
stabilize, and can the result be characterized directly from the exponent
image $J(\No)$?
\end{question}
\thref{thm:twoclosures} determines each first closure exactly. Their
incomparability makes the mixed closure a genuinely different question.
One should distinguish closure under sums from closure under all normal-form
subseries or the omega-map.

\begin{question}[Differential compatibility]
For a specified canonical surreal derivation, what is the exact locus on
which $J$ or $H_j$ commutes with it? Can a differential analogue of
\eqref{eq:explocus} be expressed using support or the iterated logarithmic
structure of the exponents?
\end{question}
The exponential locus is settled here. A derivation may depend on further
normal-form and logarithmic data, so its preservation is not a consequence
of being an ordered-field embedding. No differential claim should be
imported without verifying the actual chosen derivation.

\begin{question}[Definability in weaker signatures]
Which of $\Sk$, $\Common$, and the coefficient projection $\Pi_j$ are
definable in natural reducts of the structure with the maps $J,H_j$?
Can any of the large-cardinal information survive after the embedding
symbols and the Hadamard product are forgotten?
\end{question}
With both maps named, $\Sk$ is the immediate equalizer formula. The
question concerns losing symbols, not that trivial definition. Pure
field elementarity alone cannot recover the external strong-sum threshold.

\begin{question}[Extending the coefficient-projection isomorphism]
Under what explicit additional hypotheses does the isomorphism
$\Pi_j:J(\No)\to H_j(\No)$ extend to an automorphism of a larger
specified surreal subfield, or of the full ordered field? Which extension
requirements on monomials, simplicity, or exponentiation obstruct it?
\end{question}
The given isomorphism fixes the common field and all common monomials.
It is not globally multiplicative on $\No$, so its existing definition
cannot simply be reused as the desired extension.

\begin{question}[Forcing and lifted embeddings]
If an embedding $j$ lifts through a forcing extension, which of its old
coefficient defects are preserved and which new supports appear? Can the
lifting criterion be expressed through a commuting square of strong
companions and coefficient projections?
\end{question}
This would connect the present arithmetic construction to the repository's
cross-universe saturation work. One must first specify the ground and
extension, verify the lift, and distinguish new index sets from changed
cardinalities; no automatic preservation is asserted.

\begin{question}[A reusable formal theorem package]
Can the exact defect, image intersection, and closure theorems be formalized
in a way that separates abstract ordered Hahn algebra from the set-theoretic
embedding interface? What is the smallest verified interface sufficient
for the elementary-embedding application?
\end{question}
The proposed architecture in Appendix~\ref{app:formalization} treats the
critical-point and absoluteness obligations as explicit theorems to be
proved, not as silently trusted axioms. This would make the genuinely
new algebraic consequences independently reusable and auditable.

\section{Conclusion}
The difference between $j(I)$ and $\imset jI$ becomes a literal normal-form
difference in surreal arithmetic. Comparing the elementary-induced map
with the strongly linear map having the same monomial values exposes that
difference without sacrificing finite algebra.

The outcome is stronger than a single counterexample to summation
preservation. There is an exact equalizer, an exact intersection, a
canonical isomorphism of two incomparable immediate image fields, and an
intrinsic positive-family escape threshold. Their common field has two
different canonical enlargements: logarithmic closure produces the
elementary image, while strong-sum closure produces its Hahn companion.
One coefficient of the difference recovers the derived normal measure;
shifting the code makes it purely infinite and omnific. The same mechanism
works over complex coefficients and at supercompact seeds.

These statements answer the structural questions posed in this article
under a measurable-cardinal hypothesis. They are offered for independent
verification and comparison with the literature, not as a claim that the
large-cardinal hierarchy or a named longstanding conjecture has been
settled. The strongest next step is to test the converse: how much of the
set-theoretic embedding can be forced by the arithmetic signature alone?

\appendix
\section{Proof dependencies and a formalization plan}\label{app:formalization}

\subsection{Dependency order}
The proof does not use the repository's unreviewed set-theoretic
bi-interpretation or quotient theorems as assumptions. Its main dependency
chain is
\begin{align*}
 &\text{canonical surreal/Hahn theory}
   +\text{elementary embedding with critical point}\\
 &\hspace{12pt}\Longrightarrow
   \text{absolute normal forms and the missing-index lemma}\\
 &\hspace{12pt}\Longrightarrow J,\ H_j,
   \text{ exact defect and exponent-range exclusion}\\
 &\hspace{12pt}\Longrightarrow
   \text{equalizer, intersection, coefficient projection, strong escape}\\
 &\hspace{12pt}\Longrightarrow
   \text{omnific, surcomplex, and mask-coefficient consequences}.
\end{align*}
The logarithmic-defect homomorphism adds only the standard near-one Hahn
logarithm. The exact exponential locus and two-closure theorem add the
canonical Gonshor exponential and its purely infinite normal-form
formula. The supercompact seed theorem adds the separate
$\lambda$-sequence-closure hypothesis on $M$.

\subsection{A proposed Lean interface, not a formalization claim}
A realistic implementation should begin with a universe-bounded Hahn
workspace, then state separately how that workspace represents canonical
surreal normal forms. The following components are logically distinct.

\begin{longtable}{@{}p{0.22\textwidth}p{0.46\textwidth}p{0.25\textwidth}@{}}
\toprule
\textbf{Component} & \textbf{Obligation} & \textbf{Verification status here}\\
\midrule\endhead
Exponent reindexing & An additive order embedding preserves reverse
well-ordered support, finite convolution, and summable families. & Written
proof; repository describes related infrastructure.\\[5pt]
Canonical bridge & Sign sequences, Conway arithmetic, the omega-map, and
normal-form reconstruction agree with the Hahn representation. & Classical
input; not rebuilt or kernel-checked here.\\[5pt]
Universe absoluteness & Old normal forms and their strong sums are unchanged
on passage from $M$ to $V$. & Written reduction to the constructive
sign-expansion theorem.\\[5pt]
Embedding interface & Images of index functions; fixed coefficients;
$j(\theta)\setminus\imset j\theta$; critical-index and range-reflection
lemmas. & Written proofs; no Lean set-theory implementation supplied.\\[5pt]
Defect theorem & Split a transformed summable family into image and
missing-image indices; prove first-term and support exclusion. & Written
proof; finite deletion identity checked separately.\\[5pt]
Image geometry & Exact intersection, projection isomorphism, real closedness,
immediacy, and strong-sum escape. & Written proofs only.\\[5pt]
Measure interface & Indicator masks, Hadamard characters, completeness,
normality, and supercompact seeds. & Written proofs; finite Boolean
characters checked only as a finite analogue.\\[5pt]
Exponential interface & Absolute recursive exponential, monomial values on
purely infinite inputs, and infinitesimal power-series evaluation. &
Classical input and written consequences; not formalized here.\\
\bottomrule
\end{longtable}

The measurable-cardinal assumption must not be inserted as an unnamed
``surreal summation axiom.'' A useful interface can abstract the exact
properties of $j$ for the algebraic proofs, but the application to a normal
ultrapower still needs an explicit model of the set-theoretic construction.
That boundary is essential for avoiding a formal proof whose interesting
mathematics has merely been assumed.

\section{Verification, limitations, and an audit checklist}\label{app:verification}

\subsection{Finite regression checks}
The accompanying program \texttt{code/finite\_regression.py} uses only the
Python standard library and exact rational arithmetic. It tests finite
Laurent-polynomial reindexing, the deletion-of-indices coefficient identity,
the twisted product identity, a finite exponent shift, and the exhaustive
Boolean-algebra characters of a four-point set. Results are saved in
\texttt{data/finite\_regression\_results.json}.

These tests can detect an incorrect sign, a wrong coefficient position,
or confusion between pointwise and convolution products in a finite
identity. They do not model a nontrivial elementary embedding of the
universe. On a finite Boolean algebra every ultrafilter is principal;
the exhaustive check confirms that fact and cannot certify the existence
of a nonprincipal complete ultrafilter. There is no numerical experiment
for measurability in this package.

The PDF is compiled from the standalone source and visually inspected.
No Lean kernel, Wolfram kernel, or other proof assistant was used to
verify this article's new theorems. The source and finite program are
provided so that those limited forms of verification can be reproduced.

\subsection{Points requiring particular attention in an independent review}
The most important review target is the passage from the internal normal
form of $j(x)$ to its external canonical normal form. That passage uses
\leref{lem:absolute} and the constructive, not arbitrary, Hahn
identification. Next, verify that the missing-image exponents cannot lie
in $J(\No)$; this range-reflection argument is the source of both the
intersection theorem and the intrinsic escape theorem.

Further checks concern the separation of outer support from birthday;
the restriction to positive nonzero terms in the exact family-cardinality
criterion; the distinction between a zero coefficient probe and a zero
whole defect; and the use of Hadamard multiplication in the character
theorem. In the exponential section, check the uniform recursive
absoluteness and the monomial value of a purely infinite exponential.
The supercompact theorem requires its actual seed to belong to $M$, not
just an external set of the right apparent cardinality.

Every assertion about classes in this article reduces to a specified class
map and set-indexed operations, or to a universal property of such maps.
No proof uses a proper-class normal-form support. No claim is made that
$H_j(\No)\subseteq\No^M$, or that $J(\No)=\No^M$; neither identification
is part of the construction. All image and intersection statements occur
inside the fixed ambient $\No^V$.

\subsection{Claim summary}
Under \eqref{eq:setting}, the exact defect, support equalizer, image
intersection, projection isomorphism, positive escape, measure coefficients,
omnific and surcomplex transfers, exponential locus, and two closures are
proved in the text, relative to the stated classical inputs. The
supercompact construction is proved with its additional closure hypothesis.
The mask criterion is a proved reformulation of the classical measure
characterization. Publication priority, consistency-strength converses,
joint algebraic independence, differential compatibility, and all twelve
research questions remain outside the claims of this article.

\begingroup
\small
\raggedright
\begin{thebibliography}{99}
\setlength{\itemsep}{4pt}
\setlength{\parskip}{0pt}
\addcontentsline{toc}{section}{References}

\bibitem{Conway}
John H. Conway,
\emph{On Numbers and Games}, second edition,
A K Peters, 2001; first edition, Academic Press, 1976.
Classical source for surreal arithmetic, normal forms, and omnific integers.

\bibitem{Gonshor}
Harry Gonshor,
\emph{An Introduction to the Theory of Surreal Numbers},
London Mathematical Society Lecture Note Series 124,
Cambridge University Press, 1986.
Especially Chapters 5 (normal forms) and 10 (exponentiation).

\bibitem{vdDE}
Lou van den Dries and Philip Ehrlich,
``Fields of surreal numbers and exponentiation,''
\emph{Fundamenta Mathematicae} \textbf{167}(2) (2001), 173--188.
\href{https://doi.org/10.4064/fm167-2-3}{doi:10.4064/fm167-2-3}.
Publisher-hosted text consulted:
\url{https://www.impan.pl/shop/en/publication/transaction/download/product/89226}.

\bibitem{EhrlichNames}
Philip Ehrlich,
``Conway names, the simplicity hierarchy and the surreal number tree,''
\emph{Journal of Logic \& Analysis} \textbf{3}:1 (2011).
\url{https://logicandanalysis.org/index.php/jla/article/download/88/34}.

\bibitem{KKS}
Elliot Kaplan, Lothar Sebastian Krapp, and Michele Serra,
``Decomposing the automorphism group of the surreal numbers,''
\texttt{arXiv:2509.22374}, version 3, consulted 23 September 2026.
\url{https://arxiv.org/html/2509.22374v3}.
Used for the contemporary Hahn and strong-map framework, not as a
premise for a particular non-strong automorphism example.

\bibitem{Neeman}
Itay Neeman,
``Ultrafilters and Large Cardinals,'' author-hosted manuscript,
Section 1, consulted 23 September 2026.
\url{https://www.math.ucla.edu/~ineeman/uflc.pdf/}.
Includes the standard elementary-embedding, derived-measure, and
supercompact-seed constructions.

\bibitem{Bergman}
George M. Bergman,
``Homomorphisms on infinite direct products of groups, rings and monoids,''
\texttt{arXiv:1406.1932}, version 2, consulted 23 September 2026.
\url{https://arxiv.org/html/1406.1932v2}.
Context for ultrafilters and homomorphisms of infinite product algebras.

\bibitem{Repo}
Vladimir Reshetnikov and contributors,
\emph{Surreal}, repository snapshot
\texttt{0865f043aec113c14c69ef45006bbc7546a4e75a},
root \texttt{README.md} and \texttt{docs/README.md},
consulted 23 September 2026.
\url{https://github.com/VladimirReshetnikov/Surreal/tree/0865f043aec113c14c69ef45006bbc7546a4e75a}.
Repository statements and formalization coverage are distinguished from
independently checked proofs in the present article.

\bibitem{RepoUniverses}
\emph{Surreal Fields Across Set-Theoretic Universes},
repository research report, README and source preamble inspected at the
snapshot in \cite{Repo}.
\url{https://github.com/VladimirReshetnikov/Surreal/tree/0865f043aec113c14c69ef45006bbc7546a4e75a/docs/foundations-and-computation/surreal-fields-across-universes}.
Research context only; its new theorems are not premises here.

\bibitem{RepoBirthday}
\emph{Birthday Cutoffs and Hereditary Sets},
repository research report, README inspected at the snapshot in \cite{Repo}.
\url{https://github.com/VladimirReshetnikov/Surreal/tree/0865f043aec113c14c69ef45006bbc7546a4e75a/docs/foundations-and-computation/birthday-cutoffs-and-hereditary-sets}.
The README credits the full-class bi-interpretation announcement to
Chen--Hamkins--Yang. That announcement is not claimed as new here.

\bibitem{RepoQuotients}
\emph{Set-Sized Quotients of Omnific Integers},
catalogue entry in \texttt{docs/README.md} at the snapshot in \cite{Repo}.
Relevant report directory:
\url{https://github.com/VladimirReshetnikov/Surreal/tree/0865f043aec113c14c69ef45006bbc7546a4e75a/docs/surreal/set-sized-quotients-of-omnific-integers}.
Cited only to identify the repository's existing constant-term theme.

\end{thebibliography}
\endgroup
\end{document}
